\documentclass[11pt]{article}

\usepackage{amsmath,amssymb,amsthm}
\usepackage{geometry}
\usepackage{microtype}
\usepackage{setspace}
\usepackage{fontspec}
\usepackage{booktabs}
\usepackage{array}
\usepackage{listings}
\usepackage{xcolor}
\usepackage{hyperref}
\usepackage[nameinlink]{cleveref}

\geometry{margin=1.3in}
\setmainfont{Linux Libertine O}
\setmonofont{Linux Libertine Mono O}
\onehalfspacing

\newcommand{\ambibf}{%
  \texttt{/}\textnormal{æmbiːɛf}\texttt{/}%
}

\definecolor{codebg}{gray}{0.96}
\definecolor{codecomment}{gray}{0.45}
\definecolor{codestring}{RGB}{0,100,0}
\definecolor{codekeyword}{RGB}{0,0,160}
\definecolor{linkblue}{RGB}{30,60,140}

\lstset{
  language=Python,
  backgroundcolor=\color{codebg},
  basicstyle=\small\ttfamily,
  keywordstyle=\color{codekeyword},
  stringstyle=\color{codestring},
  commentstyle=\color{codecomment}\itshape,
  breaklines=true,
  frame=single,
  framesep=4pt,
  xleftmargin=6pt,
  xrightmargin=6pt,
  showstringspaces=false,
  tabsize=4,
}

\newtheorem{theorem}{Theorem}[section]
\newtheorem{proposition}[theorem]{Proposition}
\newtheorem{lemma}[theorem]{Lemma}
\newtheorem{corollary}[theorem]{Corollary}
\newtheorem{definition}[theorem]{Definition}
\newtheorem{remark}[theorem]{Remark}

\hypersetup{
  colorlinks=true,
  linkcolor=linkblue,
  urlcolor=linkblue,
  citecolor=linkblue,
}

\title{{\large\bfseries Globally Persistent Recomposable Thought}\\[4pt]
{\normalsize GitHub, the Intelligence Explosion, and\\
the Admissibility Geometry of Collective Cognition}}
\author{Flyxion}
\date{\today}

\begin{document}
\maketitle

\begin{abstract}
The conventional narrative attributes the current acceleration in technological
capability to the emergence of generative artificial intelligence. This essay
argues that this attribution is substantially mislocated. The more fundamental
cause is a structural transformation in the organization of human technical
cognition that preceded large language models by nearly two decades: the
creation, through GitHub and adjacent open-source infrastructure, of a globally
persistent, recursively recomposable substrate for executable thought. Generative
AI functions as a compression interface and navigational accelerant over that
substrate, not as its source. We develop this argument in five stages: first, as
a historical and infrastructural claim about the bottlenecks that GitHub
dissolved; second, as an epistemic economy argument about the incentive
structures GitHub created across academia, industry, and independent research;
third, through a concrete worked example in which Spherepop semantics are
implemented over the probabilistic substrate language \ambibf{},
demonstrating the three-tier taxonomy of admissible, boundary, and inert
repository contributions; fourth, as a connection to the Quantum SpherePop
field-theoretic framework, in which developmental geometry is interpreted as
a semantic quantization sequence and repository replay is formalized as
coarse-grained semantic reconstruction; and fifth, as a connection to the RSVP
admissibility framework, in which the GitHub substrate is understood as a
planetary-scale coherence field whose dynamics are governed by the same
lamphron--lamphrodyne duality that governs physical structure formation.
Mathematical appendices formalize the stochastic tape geometry, bubble energy,
mortality functional, stabilization operators, and attractor semantics of the
worked example.
\end{abstract}

\newpage
\tableofcontents
\newpage

%% ══════════════════════════════════════════════════════════════════════════════
\section{Introduction: Locating the Explosion}
%% ══════════════════════════════════════════════════════════════════════════════

When observers describe an intelligence explosion in progress, they typically
point at the visible outputs of large language models: the apparent fluency,
the breadth of domain competence, the speed of generation. The inference drawn
is that something qualitatively new has appeared inside the models themselves.
This essay disputes that inference --- not in order to diminish what generative
AI accomplishes, but in order to locate the causal structure of the acceleration
more precisely.

The acceleration has at least two distinct components. One is the improvement
in pattern synthesis achieved by scaling transformer architectures over massive
corpora. The other, older, and arguably more fundamental component is the
creation of a globally persistent, recomposable, machine-readable substrate of
executable human cognition. That substrate did not arrive with GPT. It arrived
with Git and GitHub, and it matured over a decade before foundation models
became publicly visible.

To conflate these two components --- to attribute the explosion to the synthesis
engine while ignoring the substrate it synthesizes over --- is to mistake the
accelerant for the fuel. The fuel is the open-source graph. GitHub
operationalized globally recomposable thought years before any language model
was capable of navigating it fluently. What generative AI contributed was a
compression interface: a mechanism for making the accumulated graph more
navigable, more searchable, more autocomplete-able. The graph itself was already
there.

Underlying this argument is an implicit redefinition of intelligence that the
essay will make progressively more explicit. Intelligence is not merely
inference efficiency within an isolated agent~\cite{russell2010}. It is the
persistence, replayability, and recomposability of developmental trajectories
across a collective substrate. GitHub is intelligent infrastructure not because
it reasons, but because it preserves and recombines the intermediate cognitive
states of millions of contributors in a form that remains recoverable,
extensible, and semantically legible across time. Generative AI is impressive
because it navigates that infrastructure fluently --- but the infrastructure is
the prior achievement.

\subsection*{A Formal Distinction}

Let $\mathcal{G}$ denote the open-source repository graph at time $t$, with
nodes $R_i$ (repositories) and directed edges representing dependency or
derivation relations. Let $\mathcal{M}_t$ denote a language model trained on
the corpus induced by $\mathcal{G}_t$. The standard narrative implicitly models
capability as a function of model parameters alone:
\[
\mathrm{capability}(t) \approx f(\mathcal{M}_t).
\]
The thesis of this essay is that this factorization is incorrect. The correct
model is:
\[
\mathrm{capability}(t) \approx f(\mathcal{M}_t,\, \mathcal{G}_t),
\]
where $\mathcal{G}_t$ is not merely training data but the admissibility geometry
over which $\mathcal{M}_t$ operates. Holding $\mathcal{G}_t$ fixed at its
pre-2005 state while scaling $\mathcal{M}_t$ to current parameter counts would
produce a system of dramatically lower practical capability, because the
substrate carrying executable procedural knowledge would be absent. The
explosion is located in the joint growth of both terms, not in $\mathcal{M}_t$
alone.

%% ══════════════════════════════════════════════════════════════════════════════
\section{What GitHub Changed Structurally}
%% ══════════════════════════════════════════════════════════════════════════════

Before distributed version control became socially normalized through GitHub,
programming knowledge was fragmented across institutional silos, mailing list
archives, FTP mirrors, academic preprints, and private internal repositories.
Code reuse existed, but composability was expensive in a deep sense: discovering
compatible dependencies, understanding undocumented build assumptions, patching
incompatibilities across codebases, and maintaining local forks all imposed
friction costs that compounded at every layer of abstraction. Knowledge transfer
was slow, lossy, and structurally conservative --- it favored large institutions
with internal memory over distributed coordination.

Git itself introduced distributed version control~\cite{git2005,torvalds2001},
but GitHub introduced something architecturally distinct: it transformed
repositories into globally addressable semantic modules. A repository ceased to
be mere storage and became instead a live executable idea, carrying versioned
history, public forking infrastructure, discoverable architecture, compositional
dependency graphs, executable documentation, collaborative debugging, social
reputation mechanisms, and machine-readable structural metadata simultaneously.

This matters because intelligence amplification in any system --- biological,
institutional, or technological --- tends to come less from increases in raw
processing capacity than from reductions in the friction of
recombination~\cite{simon1962,wiener1948}. The critical resource is not the
power of individual nodes but the composability of the graph they
inhabit~\cite{barabasi2002}.

GitHub made software development into a partially continuous evolutionary
system. A single developer anywhere on the planet can improve a tokenizer
originally written in another country, which optimizes a training pipeline
maintained in a third, which accelerates a robotics stack assembled in a fourth,
which gets incorporated into a model deployed globally within days. The causal
chain crosses borders, institutions, and time zones with a friction cost that
would have been prohibitive under any previous coordination infrastructure.
This is historically unprecedented not in degree but in kind.

The entire modern AI ecosystem is downstream of this recombination layer. The
libraries, the training frameworks, the evaluation harnesses, the inference
runtimes, the quantization methods, the adapter architectures, the diffusion
pipelines and deployment toolchains --- none of these emerged from isolated
laboratories operating independently. They emerged from an ultra-dense recursive
open-source ecosystem whose coordination was primarily mediated through GitHub.
A modern foundation model is not a singular invention. It is an enormous
stitched-together graph of repositories, and its apparent intelligence is
partially the intelligence of that graph made fluid.

In this sense GitHub functions as a planetary externalized cortex. Repositories
are memory units. Forks are evolutionary branches. Pull requests are
reconciliation operators that integrate divergent solutions. Issue trackers are
error-correction loops. Stars and watchers are salience signals that route
collective attention. Continuous integration pipelines are automated verification
mechanisms that prevent the silent accumulation of errors. The architecture
recapitulates, at planetary scale, many of the structural features of biological
working memory --- but with the crucial addition of perfect reproducibility, zero
marginal copy cost, and global simultaneous accessibility.

\subsection*{Composability as a Monotone Functional}

Define the \emph{recombination friction} $\rho(R_i, R_j)$ between two
repositories $R_i, R_j \in \mathcal{G}$ as the expected developer effort
required to integrate $R_j$ into a codebase that already depends on $R_i$.
Before GitHub, $\rho$ was high and asymmetric: it depended strongly on
institutional access, version availability, and documentation quality. GitHub
reduced $\rho$ toward a near-constant baseline by standardizing dependency
declaration, forking, and issue resolution as first-class operations.

\begin{proposition}[Composability Monotonicity]
Let $\mathcal{G}_t$ be the repository graph at time $t$ with mean pairwise
recombination friction $\bar\rho_t$. If GitHub's structural innovations reduce
$\bar\rho_t$ monotonically, then the number of reachable composite systems of
depth $d$ grows at least exponentially in $d$ as $\bar\rho_t \to 0$.
\end{proposition}

\begin{proof}[Proof sketch]
A composite system of depth $d$ requires $d$ successful integration steps. The
probability that each step succeeds is approximately $1 - \bar\rho_t / C$ for
some cost constant $C$. The expected number of reachable composites at depth
$d$ is therefore $(1 - \bar\rho_t/C)^d \cdot |\mathcal{G}_t|^d$. As
$\bar\rho_t \to 0$ this approaches $|\mathcal{G}_t|^d$, which grows
exponentially in $d$ for any fixed graph size exceeding one.
\end{proof}

\begin{corollary}
Even a modest reduction in $\bar\rho_t$ produces superlinear growth in the
accessible composite design space. GitHub's reduction was not modest; it was
structural. The consequence is the combinatorial explosion in deployable AI
infrastructure observed over the past fifteen years.
\end{corollary}

%% ══════════════════════════════════════════════════════════════════════════════

%% ══════════════════════════════════════════════════════════════════════════════
\section{Toward a Theory of Procedural Substrates}
\label{sec:procedural-substrates}
%% ══════════════════════════════════════════════════════════════════════════════

The preceding section described what GitHub changed infrastructurally. Before
proceeding to the incentive economy and the worked example, we need to establish
what kind of object a repository actually is within the framework this paper
develops. The vocabulary introduced here --- developmental trajectory,
recomposability, semantic replay, procedural substrate --- will govern all
subsequent analysis, and its stabilization early prevents ambiguity from
accumulating across later sections.

\paragraph{Artifacts versus trajectories.}
The conventional view of a software repository treats it as an information
object: a collection of files encoding a program. On this view, two repositories
with identical file contents are identical repositories, and the history of
how those contents were produced is archival metadata rather than constitutive
structure. This paper rejects that view.

A repository is not primarily an artifact. It is a compressed developmental
history --- a record of a trajectory through semantic configuration space ---
from which the artifact is merely the most recent projection. The distinction
is not merely philosophical. It determines what can be recovered from the
repository, what can be extended from it, and what training signal it
contributes to a language model. Formally:
\[
\text{artifact}(R) \;\neq\; \text{trajectory}(R).
\]
The artifact is the current state $x_n$. The trajectory is the sequence
$(x_0, x_1, \ldots, x_n)$ of states connected by the commit history. GitHub
preserves partial access to the trajectory; older publication systems preserved
only the artifact.

\paragraph{Recomposability as a geometric property.}
Recomposability is sometimes understood as a social property (whether
contributors are willing to share code) or as an infrastructural property
(whether package managers support integration). This paper uses a stricter
definition.

\begin{definition}[Recomposability]
\label{def:recomposability}
A repository $R$ is \emph{recomposable} with respect to a target architecture
$\tau$ if there exists a fork $R'$ of $R$ such that:
\begin{enumerate}
\item $R'$ satisfies the structural requirements of $\tau$,
\item the developmental trace of $R'$ is recoverable from that of $R$ plus
a bounded number of additional commits, and
\item the admissibility score $\mathcal{A}(R') \geq \mathcal{A}(R) - \varepsilon$
for some tolerance $\varepsilon > 0$.
\end{enumerate}
\end{definition}

Under this definition, recomposability is a geometric property of the
repository's position in configuration space: it asks whether $\tau$ is
reachable from $R$ via a bounded-cost trajectory that remains inside the
admissible region $\{\Phi > 0\}$. A repository may be technically open-source
but geometrically irrecomposable if its architecture is so idiosyncratic that
no bounded extension reaches a useful target. Conversely, a well-designed
modular repository may be highly recomposable even with non-trivial licensing
constraints, because its extension surfaces are clearly defined and its
dependencies are minimal.

\paragraph{Semantic replay.}
Replay is the operation by which a repository's developmental geometry is
re-instantiated in a new context. A fork is a replay operation: it takes the
developmental trace of $R$ and re-executes it under altered initial conditions
(a different base version, a different dependency set, a different target
platform). The crucial property of replay is that it need not produce an
identical result to the original: what it must produce is a result that lies
in the same attractor basin --- a system that satisfies the same
semantic invariants as the original, reached by a developmentally continuous
path.

This distinguishes replay from copying. Copying produces an identical artifact.
Replay produces a semantically equivalent trajectory endpoint, potentially via
a different route. GitHub's fork infrastructure implements replay; its snapshot
downloads implement copying. The two are not equivalent as training signals:
a corpus of copies provides no developmental geometry; a corpus of forks and
their parent repositories exposes the trajectory structure that makes
generalization possible.

\paragraph{Procedural substrates.}
We are now in a position to define the central concept of the paper.

\begin{definition}[Procedural Substrate]
\label{def:procedural-substrate}
A \emph{procedural substrate} is a globally accessible, version-controlled
repository graph $\mathcal{G}$ in which:
\begin{enumerate}
\item developmental trajectories are partially preserved (repositories have
public commit histories recoverable by any agent with network access),
\item recomposability is structurally supported (dependency declaration,
forking, and merge operations are first-class graph operations), and
\item intermediate cognitive states are externalized (issue discussions,
branch histories, and commit messages preserve reasoning traces alongside
code changes).
\end{enumerate}
The \emph{admissibility density} of a procedural substrate is
$\mathcal{A}(\mathcal{G}) = \mathbb{E}_{R \sim \mathcal{G}}[\mathcal{A}(R)]$.
\end{definition}

The thesis of this paper can now be stated precisely: the current intelligence
explosion is explained by the emergence of GitHub as the first planetary-scale
procedural substrate, and the role of generative AI is to function as a
manifold-compression operator (defined formally in \S\ref{sec:gradient})
over that substrate. Scale of compute amplifies the signal the substrate
provides; it does not substitute for the substrate itself.

\paragraph{Intermediate cognitive states.}
The third condition in Definition~\ref{def:procedural-substrate} deserves
emphasis because it represents the sharpest departure from prior publication
systems. Previous civilizations preserved \emph{finalized artifacts}: published
papers, finished products, polished monographs. GitHub preserves
\emph{intermediate procedural cognition}: the reasoning, the failure modes,
the alternatives considered, the implementation choices made and abandoned,
the debugging traces, the partial implementations, and the evolutionary
trajectory from naive first attempt to refined architecture.

This matters because intermediate states carry the developmental geometry that
final states conceal. A published paper explaining an algorithm does not expose
the five failed implementations that preceded it, the bug that forced a
redesign, or the off-by-one error that revealed a deeper architectural
assumption. A public repository with its full commit history does. The
training signal embedded in that history is qualitatively different from the
training signal in the final published artifact: it exposes not only what
the system is but how it became what it is, which is precisely the developmental
geometry that enables both human extension and model generalization.

\section{The Epistemic Economy: Academia, Industry, and the Incentive Inversion}
%% ══════════════════════════════════════════════════════════════════════════════

The structural argument becomes considerably stronger when one examines what
GitHub did simultaneously to academia and to industry --- and in particular the
strange new incentive landscape it created across both.

Traditional academic knowledge production was optimized around scarcity.
Journal slots were scarce. Institutional legitimacy was scarce. Peer review
was slow and systematically selected for finished narratives, clean positive
outcomes, and retrospective theoretical coherence. The payoff structure strongly
discouraged publishing negative results, failed experiments, intermediate
implementations, and unfinished tooling. Much of the productive cognitive work
behind a published result disappeared permanently into private hard drives and
unpublished notes.

Industry was similarly siloed. Tooling was proprietary, research pipelines were
closed, engineering cultures were isolated, and the primary mechanism for
knowledge transfer across institutions was credential-mediated hiring rather than
artifact-mediated accumulation.

GitHub partially inverted both systems simultaneously.

For researchers, a public repository became a mechanism for reproducibility and
a new form of epistemic credibility. A paper without accompanying code has come
to look increasingly incomplete in any field where results depend on exact
preprocessing pipelines, training configurations, dependency version constraints,
hidden implementation choices, and evaluation harnesses that are not fully
specifiable in prose. The repository functions as an executable appendix: not
merely supplementary material, but primary evidence. The claim shifts from the
assertion of a result to the demonstration of the exact machinery that generated
it. This is a profound epistemological change, because it substantially narrows
the gap between claim and verifiable warrant.

More unexpectedly, GitHub rewards incompleteness in ways that the traditional
academic publication system strongly punished. A half-working repository can
generate enormous downstream value if someone else forks it, repairs a critical
bug, integrates it with another project, reinterprets its architecture, or
extracts a single useful subroutine that would otherwise have remained invisible.
A failed transformer experiment may contain a novel architectural idea. An
abandoned rendering engine may contain a useful data structure. An incomplete
theorem prover may implement a technique that another researcher needs. Under
the old system, all of this would have vanished. GitHub externalizes cognitive
debris instead of discarding it, and evolutionary systems require variation ---
including failed variation --- not merely success. This transforms failure itself
into reusable infrastructure, which is a genuinely novel property of the
ecosystem GitHub produced.

Businesses reinforce this dynamic because repositories also became the
highest-bandwidth talent discovery mechanism available~\cite{benkler2006,raymond1999}.
A repository is substantially more informative than a r\'esum\'e. Code quality,
architectural taste, debugging persistence, documentation discipline,
collaboration behavior, design philosophy, and the trajectory of intellectual
growth over time are all directly visible in a sufficiently active public
profile. This creates a strange new incentive structure: people publish not only
because publication is epistemically virtuous but because publication is
economically actionable. The repository is simultaneously a research artifact,
a portfolio, a social signal, a collaboration node, a recruiting surface, and a
contribution to the commons. All of these payoffs align toward the same
behavior: public technical disclosure.

Even large corporations participate in this economy because selective open
sourcing can attract talent, establish de facto standards, shape downstream
ecosystems, reduce internal maintenance costs, appropriate community labor, and
increase platform lock-in. The incentive structures across academic researchers,
independent hobbyists, early-stage startups, and large technology corporations
are therefore partially aligned toward an outcome --- global public technical
disclosure --- that no centralized actor planned and no traditional institution
could have organized.

That alignment is historically unusual. In previous eras, academics hoarded
unfinished work, corporations concealed infrastructure, hobbyists lacked
visibility, and failed ideas vanished. GitHub changed the payoff matrix at every
level of the ecosystem simultaneously, and the result was an unprecedented
externalization of procedural knowledge into globally accessible executable form.

\subsection*{The Payoff Matrix Inversion}

Let $V_\mathrm{pub}(R)$ denote the private value to a developer of publishing
repository $R$, and $V_\mathrm{priv}(R)$ the value of keeping it private.
Under the pre-GitHub regime, $V_\mathrm{priv}(R) > V_\mathrm{pub}(R)$ for
most $R$: secrecy preserved competitive advantage, and publication infrastructure
was costly. GitHub altered the components of $V_\mathrm{pub}$ by adding
reputation signals $\sigma(R)$, talent-market visibility $\tau(R)$, and
community maintenance value $\mu(R)$:
\[
V_\mathrm{pub}(R) = V_0(R) + \sigma(R) + \tau(R) + \mu(R),
\]
where $V_0(R)$ is the baseline epistemic value of the artifact and $\sigma,
\tau, \mu \geq 0$. For repositories of sufficiently broad utility,
$\sigma + \tau + \mu$ dominates any competitive-secrecy premium, reversing the
inequality. The result is a Nash equilibrium shift: once enough participants
publish, the reputational penalty for not publishing exceeds the competitive
benefit of secrecy, making publication the dominant strategy across the
ecosystem.

\begin{proposition}[Publication Dominance]
In a market where talent acquisition depends on public repository visibility and
where standard-setting requires ecosystem adoption, there exists a threshold
$\theta > 0$ such that for any repository with aggregate utility exceeding
$\theta$, publication is the dominant strategy regardless of initial competitive
considerations.
\end{proposition}

This threshold is endogenous: as more participants publish, $\tau(R)$ increases
for all repositories (because the talent market becomes more repository-literate)
and $\sigma(R)$ increases (because visibility scales with network density),
lowering $\theta$ further. The equilibrium is therefore self-reinforcing once
critical mass is reached.

%% ══════════════════════════════════════════════════════════════════════════════
\section{The Recursion That Closed}
%% ══════════════════════════════════════════════════════════════════════════════

The current acceleration is therefore the product of a closed recursion loop
rather than a single technological breakthrough. GitHub accumulated global
executable cognition over roughly a decade and a half. Large language models
trained on that cognition. Those models now assist in generating new
repositories. Those repositories feed future training corpora. The cycle
accelerates.

The explosion is not located inside the model alone. It is located in the
coupling between globally persistent open-source infrastructure, distributed
version control, internet-scale collaborative coordination, and machine-assisted
synthesis. Remove any of these components and the feedback loop either slows
dramatically or ceases.

One useful formulation distinguishes the memory layer from the synthesis layer.
GitHub created civilization-scale working memory: a globally accessible,
perfectly reproducible, continuously updated record of executable technical
cognition accumulated across millions of contributors over years. Generative AI
created civilization-scale autocomplete over that memory: a mechanism for
navigating, compressing, translating, and recombining the contents of that record
with substantially reduced friction. The autocomplete is impressive. But it
autocompletes over a substrate it did not create.

An LLM trained exclusively on text that predated the open-source ecosystem would
be dramatically weaker in practical utility, because the open-source ecosystem is
not merely training data --- it is the ontological substrate from which the
model's representations of procedural reasoning are derived. GitHub may not
merely have enabled AI. It may have manufactured the very conceptual vocabulary
that AI learned to navigate.

The deeper bottleneck that was dissolved was never insufficient raw intelligence.
It was insufficient coordination, reuse, interoperability, discoverability, and
cumulative persistence of technical cognition across institutional and national
boundaries. GitHub solved that bottleneck first. Generative AI subsequently
accelerated throughput over a substrate that had already been built.

\subsection*{Fixed-Point Analysis of the Recursion}

Let $\mathcal{G}_t$ denote the repository graph and $\mathcal{M}_t$ a language
model at time $t$. Define the joint state $\Sigma_t = (\mathcal{G}_t,
\mathcal{M}_t)$. The recursion operates through two coupled update rules:
\begin{align}
\mathcal{M}_{t+1} &= \mathrm{Train}(\mathcal{M}_t,\; \mathcal{G}_t), \label{eq:model-update}\\
\mathcal{G}_{t+1} &= \mathrm{Extend}(\mathcal{G}_t,\; \mathcal{M}_t), \label{eq:graph-update}
\end{align}
where $\mathrm{Train}$ updates model parameters from the current graph corpus
and $\mathrm{Extend}$ adds repositories generated or assisted by the current
model. The composition $\Phi = \mathrm{Extend} \circ \mathrm{Train}$ defines a
dynamical system on the joint space.

\begin{proposition}[Recursion Amplification]
If $\mathrm{Train}$ is monotone in graph quality (better-structured graphs
produce more capable models) and $\mathrm{Extend}$ is monotone in model
capability (more capable models generate higher-quality repositories), then the
joint dynamical system is monotone increasing in both coordinates simultaneously,
and no stable fixed point exists below the resource ceiling of the ecosystem.
\end{proposition}

\begin{proof}[Proof sketch]
Let $q(\mathcal{G})$ denote graph quality and $c(\mathcal{M})$ model capability,
both real-valued. By assumption $c(\mathcal{M}_{t+1}) \geq c(\mathcal{M}_t)$
whenever $q(\mathcal{G}_t) \geq q(\mathcal{G}_{t-1})$, and $q(\mathcal{G}_{t+1})
\geq q(\mathcal{G}_t)$ whenever $c(\mathcal{M}_t) \geq c(\mathcal{M}_{t-1})$.
By induction on $t$, both sequences are non-decreasing. A fixed point
$\Sigma^*$ would require $\mathrm{Train}(\mathcal{M}^*, \mathcal{G}^*) =
\mathcal{M}^*$ and $\mathrm{Extend}(\mathcal{G}^*, \mathcal{M}^*) =
\mathcal{G}^*$ simultaneously, which cannot hold below the resource ceiling
because each update strictly improves upon the previous state.
\end{proof}

\begin{corollary}
The feedback loop described in equations \eqref{eq:model-update} and
\eqref{eq:graph-update} is self-accelerating. The observed nonlinearity in
AI capability growth is not a property of the model architecture alone; it
is a property of the coupled dynamical system.
\end{corollary}

%% ══════════════════════════════════════════════════════════════════════════════

%% ══════════════════════════════════════════════════════════════════════════════
\section{Why Esolangs Matter Epistemically}
%% ══════════════════════════════════════════════════════════════════════════════

Before presenting the worked example, it is worth explaining why an esoteric
programming language --- rather than a more conventional computational model ---
provides the most instructive substrate for the argument. The choice is not
arbitrary or ornamental. Esolangs function as compressed laboratories for
admissibility geometry because they isolate computational invariants under
adversarial conditions.

A conventional programming language is designed to minimize friction between
programmer intent and machine execution. It supplies deterministic control flow,
precise arithmetic, predictable memory, and a type system that enforces
consistency at the symbolic level. These features are engineered precisely to
eliminate the gap between syntactic intent and operational behavior. This makes
conventional languages excellent tools for building systems, but poor
instruments for studying the minimal conditions under which computation
survives substrate unreliability.

An esoteric language designed adversarially --- one that deliberately removes
the scaffolding that conventional languages supply --- forces exactly this
question into the open: what is the irreducible structural minimum that permits
coherent computation? The answer reveals which invariants are genuinely load-bearing
and which are engineering conveniences.

\ambibf{} is particularly well suited to this purpose because it removes
\emph{directional determinism}: the most basic assumption of Turing-complete
computation, namely that symbolic control flow determines operational behavior,
is destroyed at the syntactic level. The programmer can write \texttt{+} or
\texttt{-}, \texttt{<} or \texttt{>}, but the runtime erases the distinction.
What survives is only what the substrate cannot destroy: in this case, parity.

\paragraph{Symbolic determinism versus statistical semantic persistence.}
This establishes the central epistemological distinction that bridges the
worked example to the larger GitHub argument. Ordinary programming languages
stabilize computation through \emph{symbolic determinism}: the semantics of
a program are fully determined by its text, and execution is a pure function
of the symbolic state. Spherepop-style systems stabilize computation through
\emph{statistical semantic persistence}: meaning is not encoded in any single
symbolic state but in the distribution of states that a system reliably
occupies under perturbation. A bubble is not a variable whose value is the
current tape cell content; it is a localized probability distribution whose
center of mass defines the semantic state.

This distinction is not merely technical. It is the bridge between:
\begin{itemize}
\item the GitHub replay argument (repository identity is defined not by a
snapshot but by the distribution of states a development process reliably
returns to under forking and perturbation),
\item the stochastic semantics of \ambibf{} and Spherepop (meaning persists
through attractor stability rather than symbolic identity),
\item the LLM training argument (models learn developmental geometry ---
the attractor structure of good code --- not merely surface syntax),
\item the RSVP admissibility framework (the scalar field $\Phi > 0$ encodes
regions of configuration space where attractor stability is maintained), and
\item open-source recomposition (composability is possible because recomposable
systems share attractor structure across their extension surfaces, not merely
compatible type signatures).
\end{itemize}

Without this bridge, the transition from Python interpreters to field-theoretic
language may seem abrupt. With it, the progression is conceptually continuous:
each layer of formalism describes the same phenomenon --- the persistence of
recoverable developmental structure under perturbation --- at a different scale
and resolution.

\paragraph{The esolang as minimal model organism.}
In this sense \ambibf{} functions as what biologists call a model organism:
a system deliberately chosen for its tractability rather than its
representativeness, from which general principles can be extracted and
transferred to less tractable domains. Just as \textit{C. elegans} reveals
principles of developmental biology that transfer to vertebrates, \ambibf{}
reveals principles of admissibility geometry that transfer to GitHub repositories,
LLM training dynamics, and RSVP field theory. The worked example is not a
detour from the main argument. It is the main argument at its most legible scale.

\section[A Worked Example]{A Worked Example: Spherepop over \ambibf{}}
\label{sec:worked}
%% ══════════════════════════════════════════════════════════════════════════════

The argument so far has remained at the level of ecosystem dynamics. To make
it concrete, we need a worked example: an actual repository, developed in
stages, whose publication history instantiates the theoretical claims rather
than merely illustrating them. The example chosen here is deliberately
non-trivial. It is not a utility script or a data-processing tool. It is an
implementation of Spherepop semantics over a probabilistic substrate language
called \ambibf{}, and its interest lies precisely in the fact that the system
it implements is itself a model of how stable meaning emerges from unreliable
substrates. The repository becomes, in this sense, a miniature of the larger
ecosystem it is embedded in.

\ambibf{} is a Brainfuck derivative created by ais523 in 2012. Its design is
deliberately adversarial toward determinism. Where Brainfuck provides distinct
increment and decrement operators, \ambibf{} collapses them into a single
instruction that randomly does either. Where Brainfuck provides distinct left
and right pointer movements, \ambibf{} collapses these into a single instruction
that randomly does either. The bracket instructions \texttt{[} and \texttt{]}
retain their Brainfuck semantics, but since the arithmetic and movement operators
are both nondeterministic, the control flow they govern operates over a
stochastic field rather than a deterministic memory. The tape is infinite in
both directions and every cell holds a bignum value, which provides unexpectedly
rich storage capacity, but that capacity is accessible only through
probabilistic navigation.

The consequence is philosophically significant. This is not merely randomized
Brainfuck. The randomness is structural: every directional and arithmetic
intention in the source code is erased at the operational level. A program that
writes \texttt{+} and a program that writes \texttt{-} are operationally
identical. The source code contains apparent semantic content that the runtime
systematically destroys. Syntax no longer cleanly determines execution. Meaning
exists only as a distribution over possible trajectories.

This raises a question that makes the language genuinely interesting as a
theoretical object rather than a mere curiosity: how much determinism is
actually necessary for computation? The Turing-completeness of \ambibf{}
remains an open question, but the analysis in the original language specification
demonstrates that reliable computation can emerge from unreliable
primitives~\cite{neumann1956}. By encoding counter values not as single tape
cells but as parity-structured fields distributed across tape regions, the
language allows certain invariants to be recovered statistically even when exact
cell values cannot be trusted. Parity remains one of the few geometric
properties that survives noisy drift with recoverable probability. This connects
\ambibf{} to von Neumann's classical result on reliable computation from
unreliable components~\cite{neumann1956}, to error-correcting codes~\cite{cover2006},
and to the broader family of systems where stable structure emerges from
stochastic substrates~\cite{wolfram2002}.

A minimal Python interpreter for \ambibf{} captures these dynamics directly:

\begin{lstlisting}
import random
from collections import defaultdict

class AmbiBFMachine:
    def __init__(self):
        self.tape = defaultdict(int)
        self.ptr = 0

    def step(self, command):
        if command in "+-":
            self.tape[self.ptr] += random.choice([-1, 1])
        elif command in "<>":
            self.ptr += random.choice([-1, 1])

    def run(self, code, max_steps=100000):
        brackets = self.match_brackets(code)
        ip = 0
        steps = 0
        while ip < len(code) and steps < max_steps:
            c = code[ip]
            if c == "[":
                if self.tape[self.ptr] == 0:
                    ip = brackets[ip]
            elif c == "]":
                if self.tape[self.ptr] != 0:
                    ip = brackets[ip]
            else:
                self.step(c)
            ip += 1
            steps += 1
        return dict(self.tape)

    @staticmethod
    def match_brackets(code):
        stack, pairs = [], {}
        for i, c in enumerate(code):
            if c == "[":
                stack.append(i)
            elif c == "]":
                if not stack:
                    raise SyntaxError("Unmatched ]")
                j = stack.pop()
                pairs[i] = j
                pairs[j] = i
        if stack:
            raise SyntaxError("Unmatched [")
        return pairs
\end{lstlisting}

This is the lamphron moment: a compact, specialized, self-contained artifact
that makes a single architectural commitment visible. The bracket-matching
function is deterministic. The tape is a \texttt{defaultdict}. The \texttt{step}
function delegates all nondeterminism to a single \texttt{random.choice} call.
Every design decision encodes an irreversible constraint. The choice to use a
hash map rather than a list encodes a commitment to infinite sparse tape. The
\texttt{max\_steps} guard encodes a commitment to termination safety. The bracket
matching as a preprocessing pass rather than a runtime scan encodes a commitment
to $O(1)$ jump cost. These are Spherepop events: irreversible constraints that
narrow the option-space of future extensions while keeping the trajectory inside
the admissible region $\Phi > 0$.

A model ingesting this repository does not memorize the file. It updates its
representation of what it looks like when a stochastic tape machine is
implemented carefully: what the bracket-matching idiom looks like, what the
\texttt{defaultdict} pattern signals about tape geometry, what the separation
between \texttt{step} and \texttt{run} reveals about the intended extension
surface. These are procedural patterns that transfer. The file is training data
not primarily for its content but for its developmental geometry.

\subsection*{Parity Invariance Under Stochastic Drift}

The key insight enabling reliable computation on this substrate is that parity
is preserved under symmetric random walks. We formalize this as follows.

\begin{lemma}[Parity Preservation]\label{lemma:parity}
Let $p(t) \in \mathbb{Z}$ be the tape pointer evolving under the \ambibf{}
movement rule: at each step, $p(t+1) = p(t) + \eta_t$ where $\eta_t \in
\{-1, +1\}$ with equal probability. Then the parity $p(t) \bmod 2$ is
deterministic given $p(0)$: $p(t) \bmod 2 = (p(0) + t) \bmod 2$ for all $t$.
\end{lemma}

\begin{proof}
Since $\eta_t \in \{-1, +1\}$, we have $\eta_t \equiv 1 \pmod{2}$ in both
cases. Therefore $p(t+1) \equiv p(t) + 1 \pmod{2}$ at every step, regardless
of the random outcome. By induction, $p(t) \equiv p(0) + t \pmod{2}$.
\end{proof}

\begin{corollary}[Dead Reckoning]
A program executing on \ambibf{} can always determine the current parity of
the tape pointer by counting steps, even with no access to the pointer's
absolute position. This is the foundational invariant on which the counter
encoding described in the original language specification rests.
\end{corollary}

\begin{remark}
The lemma shows that parity is the \emph{unique} positional invariant
preserved under arbitrary symmetric random walks on $\mathbb{Z}$: all other
modular invariants are destroyed by stochastic drift. The \ambibf{} encoding
strategy is therefore not merely clever but optimal given the substrate's
constraints.
\end{remark}

%% ══════════════════════════════════════════════════════════════════════════════
\section{Bubbles as Metastable Attractors: The Spherepop Layer}
%% ══════════════════════════════════════════════════════════════════════════════

The \texttt{AmbiBFMachine} alone is a substrate interpreter. It demonstrates
that \ambibf{} dynamics can be simulated, but it does not yet raise the question
that makes the system theoretically interesting: can coherent semantic structures
persist on this substrate despite continuous stochastic erosion? Answering this
question requires a second layer, and the second layer is Spherepop.

In the Spherepop formalism, a bubble is not a variable. It is a named localized
field whose value emerges through stabilization rather than assignment. Applied
to an \ambibf{} substrate, a bubble is a parity-preserving neighborhood of tape
cells whose collective behavior encodes a semantic state that survives repeated
perturbation. The parity constraint is essential: it is one of the few geometric
properties that can be recovered statistically after noisy pointer drift, because
the even-odd structure of tape positions is invariant under the symmetric random
walk (see Appendix~\ref{sec:bubbleregions} for the formal definition of bubble
support regions).

A minimal Spherepop layer over the interpreter captures this directly:

\begin{lstlisting}
class Spherepop:
    def __init__(self):
        self.machine = AmbiBFMachine()
        self.bubbles = {}
        self.next_cell = 0

    def bubble(self, name):
        self.bubbles[name] = self.next_cell
        self.next_cell += 2

    def focus(self, name):
        target = self.bubbles[name]
        while self.machine.ptr != target:
            if self.machine.ptr < target:
                self.machine.ptr += 1
            else:
                self.machine.ptr -= 1

    def jiggle(self, name, amount=1):
        self.focus(name)
        for _ in range(amount):
            self.machine.run("+")

    def pop(self, name):
        self.focus(name)
        self.machine.run("[+]")

    def stabilize(self, name, attempts=101):
        self.focus(name)
        values = [self.machine.tape[self.machine.ptr]
                  for _ in range(attempts)]
        positives = sum(1 for v in values if v > 0)
        negatives = sum(1 for v in values if v < 0)
        if positives > negatives:
            return "+bubble"
        elif negatives > positives:
            return "-bubble"
        else:
            return "void"

    def read(self):
        return {name: self.machine.tape[cell]
                for name, cell in self.bubbles.items()}
\end{lstlisting}

Four operations define the minimal Spherepop interface. \texttt{bubble}
allocates a named region with spacing of two cells to reduce parity
contamination between neighbors. \texttt{focus} navigates deterministically
to the bubble center --- the one place where the Spherepop layer reasserts
spatial authority over the drifting substrate. \texttt{jiggle} perturbs the
bubble stochastically, accumulating energy in the local field. \texttt{pop}
applies the \texttt{[+]} idiom, injecting symmetric perturbations until the
local field decoheres to zero: not a variable assignment, but an entropic
collapse event. \texttt{stabilize} classifies the bubble by repeated
observation, returning a semantic state determined by majority vote. This is
not reading a memory address. It is a measurement over a distribution.

The mortality consequence is immediate and formally sharp. Left to evolve under
unbiased \ambibf{} dynamics, every bubble eventually decoheres to zero with
probability one (Appendix~\ref{sec:mortality}, Theorem~\ref{thm:mortality}).
Semantic structures on this tape are mortal. They persist only through active
stabilization, which means that computation on this substrate is not the
execution of a static program but the ongoing governance of metastable attractors
against entropic drift.

The lamphrodyne extension takes the bubble metaphor seriously as a field
phenomenon rather than a naming convention:

\begin{lstlisting}
class Bubble:
    def __init__(self, center, parity, radius):
        self.center = center
        self.parity = parity
        self.radius = radius

def sample_field(machine, bubble):
    return [machine.tape[i]
            for i in range(bubble.center - bubble.radius,
                           bubble.center + bubble.radius + 1)
            if abs(i % 2) == bubble.parity]

def semantic_state(values):
    energy = sum(abs(v) for v in values)
    polarity = sum(values)
    if energy == 0:
        return "void"
    if polarity > 0:
        return "coherent-positive"
    if polarity < 0:
        return "coherent-negative"
    return "chaotic"
\end{lstlisting}

This is the lamphrodyne moment: the single-cell bubble model has been integrated
into a richer semantic architecture where meaning emerges from field
distributions, stabilization is a regional governance process, and the boundary
between coherent and chaotic states is a continuous quantity rather than a binary
flag. The two-stage development is itself a small instance of the
lamphron--lamphrodyne feedback loop that the essay has been describing at
ecosystem scale.

\subsection*{Stabilization Threshold and Majority Vote}

The \texttt{stabilize} method implements a majority vote over a sample of size
$n = 101$. We derive the minimum number of samples needed to guarantee a
reliable classification with high probability.

Let $p \in (1/2, 1]$ be the true probability that a tape cell has positive
value at any given observation time. The majority vote returns ``\texttt{+bubble}''
if more than $n/2$ observations are positive.

\begin{proposition}[Sample Sufficiency for Stabilization]
Let $X_1, \ldots, X_n$ be independent Bernoulli$(p)$ observations of a bubble
cell sign ($1$ for positive, $0$ otherwise). The majority vote estimator
$\hat{p} = \mathbf{1}[\sum_i X_i > n/2]$ satisfies
\[
\Pr\bigl(\hat{p} \neq \mathbf{1}[p > 1/2]\bigr)
\;\leq\; 2\exp\!\Bigl(-2n(p - 1/2)^2\Bigr)
\]
by Hoeffding's inequality. For $p - 1/2 = \delta$ and error tolerance
$\varepsilon > 0$, the sufficient sample size is
\[
n \;\geq\; \frac{\ln(2/\varepsilon)}{2\delta^2}.
\]
\end{proposition}

\begin{corollary}
For a bubble with coherence advantage $\delta = 0.1$ (ten percent more positive
observations than negative) and error tolerance $\varepsilon = 0.01$, the
required sample size is $n \geq \lceil \ln(200)/0.02 \rceil = 266$. The
implementation uses $n = 101$, which is sufficient for $\delta \geq 0.15$
at the same confidence level. A production implementation should parameterize
$n$ as a function of the expected substrate noise level.
\end{corollary}

%% ══════════════════════════════════════════════════════════════════════════════
\section{The Three Tiers of Admissibility}
%% ══════════════════════════════════════════════════════════════════════════════

The important property of the Spherepop interpreter is not that it works
perfectly, but that its failures remain geometrically legible. Even when
stabilization collapses --- when the governance coefficient is too small, when
the bubble radius is poorly chosen, when the \texttt{focus} function navigates
to the wrong parity class --- the system retains a recoverable semantic
topology. The parity structure, the replay logic, the attractor regions, and
the mortality-management strategies remain inferable from the event history. A
reader encountering a broken version of this interpreter can reconstruct what
kind of system it was trying to become. This distinction between recoverable
semantic failure and irrecoverable noise is precisely what separates repositories
that contribute to the collective procedural substrate from uploads that remain
informationally inert.

To make this distinction rigorous, we must separate two quantities that are
frequently conflated in discussions of training data quality: informational
content in the Shannon sense~\cite{shannon1948}, and admissible replay structure
in the Spherepop sense. These quantities are orthogonal. A random byte stream
can carry arbitrarily high Shannon entropy while containing zero recoverable
procedural trajectory. Conversely, a partially broken repository may contain
syntax errors, dead code, and incomplete implementations while still exposing
extremely rich developmental geometry. The crucial quantity for training manifold
contribution is not raw information density but admissible replay structure: the
degree to which an artifact's event history can be reconstructed and extended by
a reader --- human or model --- encountering it.

This gives a three-tier taxonomy of repository contributions, grounded in the
admissibility geometry of the preceding sections.

The first tier comprises well-formed repositories with stable semantic normal
forms. The Spherepop interpreter, published with documentation and mathematical
appendices, is an example. Every design decision encodes an irreversible
constraint that narrows the option-space of future extensions while maintaining
the trajectory inside $\Phi > 0$. A model ingesting this repository learns the
parity structure of the tape, the stabilization loop pattern, the relationship
between bubble energy and semantic persistence, the Python idioms for
probabilistic field dynamics, and the mathematical vocabulary connecting entropy
to coherence. It learns these things not by memorizing the file but by updating
its representation of what coherent procedural reasoning over stochastic
substrates looks like across multiple levels of abstraction simultaneously.

The second tier comprises broken repositories with recoverable architectural
intent. Consider a version of the Spherepop interpreter in which the
\texttt{stabilize} method contains an off-by-one error in the sample window,
causing systematic bias in bubble classification at boundary conditions. The
interpreter is wrong, but it still encodes the parity assumption, the
majority-vote protocol, the three-way distinction between positive, negative,
and void semantic states, and the architectural decision to separate measurement
from governance. A fork that corrects the off-by-one error leaves all of these
intact. In the RSVP framework this repository sits near $Z(\Phi)$: on the
admissibility boundary, potentially crossable via forking.

The third tier comprises uploads with no recoverable semantic structure. A table
of random numbers occupies a categorically different region from either of the
first two tiers. It contains no Spherepop normal form because there is no
sequence of irreversible design decisions to reconstruct. There is no parity
assumption to infer, no stabilization strategy to extend, no attractor geometry
to fork. A recording of radio static is similarly outside the admissibility
manifold: its Shannon entropy may be high, but its developmental vector field
is null. No fork of a static recording becomes a Spherepop interpreter.

A repository's contribution to the collective substrate is therefore proportional
not to its size, not to its Shannon entropy, and not even to its correctness,
but to the recoverability of its semantic attractor under perturbation:
\[
\mathcal{A}(R)
\;\propto\;
\frac{\text{recoverable semantic structure}}{\text{perturbation magnitude}}.
\]

\subsection*{Orthogonality of Shannon Entropy and Replay Structure}

\begin{proposition}[Independence of Information Measures]
Shannon entropy $H$ and admissible replay structure $\mathcal{A}$ are
statistically independent in the space of all byte sequences. That is, for any
values $h$ and $a$, there exist repositories with $H = h$ and $\mathcal{A} = a$
in all four combinations of high/low.
\end{proposition}

\begin{proof}[Proof sketch]
High $H$, high $\mathcal{A}$: a large well-formed codebase with many distinct
tokens. Low $H$, high $\mathcal{A}$: a minimal interpreter using a small
alphabet (e.g., the \texttt{AmbiBFMachine} itself). High $H$, low $\mathcal{A}$:
a table of uniform random bytes --- maximal Shannon entropy, zero architectural
intent. Low $H$, low $\mathcal{A}$: a file of identical bytes --- minimal
entropy, no recoverable design trajectory. All four cases are constructible,
so the measures are not monotonically related in either direction.
\end{proof}

\begin{corollary}
Filtering training corpora by Shannon entropy (a standard quality heuristic)
does not reliably select for high $\mathcal{A}$. A corpus filter optimized
for developmental geometry would require a distinct criterion: something closer
to compressibility of the commit history, architectural coherence of the
dependency graph, or forkability under perturbation.
\end{corollary}

%% ══════════════════════════════════════════════════════════════════════════════

%% ══════════════════════════════════════════════════════════════════════════════
\section{Commit Histories as Temporal Geometry}
%% ══════════════════════════════════════════════════════════════════════════════

The three-tier admissibility taxonomy classifies repositories by the
recoverability of their semantic attractor. But repositories are not static
objects: they are processes unfolding in time, and their admissibility structure
is encoded not only in their current state but in the discretized trace through
semantic configuration space that their commit history records. This section
formalizes that observation and connects it to the Lamport clock structure
underlying distributed version control.

\paragraph{Commits as irreversible constraint applications.}
Each commit to a repository applies a transformation $T_n$ to the current
semantic state $x_n$, producing $x_{n+1} = T_n(x_n)$. In the RSVP framework,
this corresponds to an increment in the entropy field $S$ along the development
trajectory: each $T_n$ that moves the codebase further into $\{\Phi > 0\}$
represents an irreversible narrowing of the future option-space. The commit
sequence $(T_1, T_2, \ldots, T_n)$ is therefore the Spherepop normal form of
the repository: the ordered sequence of irreversible decisions whose composition
constitutes the repository's semantic identity.

\begin{definition}[Developmental Trace]
The \emph{developmental trace} of a repository $R$ with commit history
$(T_1, \ldots, T_n)$ is the sequence of states
\[
x_0, x_1 = T_1(x_0), x_2 = T_2(x_1), \ldots, x_n = T_n(x_{n-1}),
\]
together with the associated constraint tightening sequence
$S_0 \leq S_1 \leq \cdots \leq S_n$ induced by the entropy field. The trace
is \emph{recoverable} if the sequence $(T_1, \ldots, T_n)$ can be
reconstructed from partial evidence with bounded error.
\end{definition}

A Tier~1 repository has a fully recoverable developmental trace: every design
decision is inferrable from the public commit history and accompanying
documentation. A Tier~2 repository has a partially recoverable trace: some
commits are missing, some transformations are ambiguous, but the overall
architectural trajectory remains inferable. A Tier~3 upload has no
recoverable trace: there is no sequence of intentional transformations to
reconstruct.

\paragraph{Commit hashes as Lamport timestamps.}
This connects to a precise analogy raised by the Lamport clock literature. A
Lamport timestamp~\cite{lamport1978} assigns a logical time $q_t$ to each
event in a distributed system such that causal precedence is preserved: if
event $A$ causally precedes event $B$, then $q_A < q_B$. Critically, Lamport
timestamps encode causal order without requiring global synchronization --- they
are the minimal invariant that survives distributed, asynchronous computation.

Git commit hashes are the executable analogue of Lamport timestamps. Each
commit hash encodes the complete causal history of the repository up to that
point: it is a cryptographic summary of all prior commits, all file contents,
and the commit message. Like a Lamport clock, the hash encodes causal order
without requiring central coordination. Like parity in \ambibf{}, it is the
minimal invariant that survives the distributed, asynchronous process of
collaborative development.

\begin{proposition}[Commit Hashes as Causal Parity]
Let $q_t$ denote the commit hash at step $t$ in a repository's history. Then:
\begin{enumerate}
\item $q_t$ determines a unique causal past for the repository at time $t$
(by the collision resistance of the underlying hash function);
\item if $q_t$ is known, the developmental trace $(T_1, \ldots, T_t)$ is
in principle recoverable from the Git object store;
\item two repositories with the same $q_t$ are semantically identical at
time $t$ regardless of when, where, or by whom they were created.
\end{enumerate}
Commit hashes are therefore the distributed version control analogue of
parity bits: the unique causal invariant that survives stochastic, asynchronous
development and permits dead reckoning of developmental position.
\end{proposition}

The analogy extends further. Just as parity preservation in \ambibf{} is
the foundational invariant enabling metastable computation on a noisy
substrate (Lemma~\ref{lemma:parity}), commit-hash integrity is the foundational
invariant enabling semantic persistence in the distributed, asynchronous
development process. A repository without intact commit history --- a snapshot
without provenance --- loses its developmental trace in exactly the way a
bubble loses its parity structure when the pointer drifts to an unknown
position. The admissibility of the artifact depends on the recoverability
of its causal history, not only on the content of its current state.

\paragraph{Issue threads and branch histories as developmental debris.}
Beyond commits, repositories typically expose additional layers of developmental
geometry: issue threads record the reasoning behind architectural decisions and
the failure modes encountered during development; branch histories record
exploratory trajectories that were not merged; pull request discussions record
the deliberation process by which competing extensions were adjudicated. All of
these are instances of what we have called intermediate cognitive states ---
externally persistent records of the developmental process that would, under
any previous publication system, have been discarded as ancillary metadata.

In the present framework they are not ancillary. They are the primary evidence
from which the developmental trace is reconstructible. A repository whose
issue threads document a debate about whether to use majority-vote or
maximum-likelihood stabilization exposes more developmental geometry than one
whose commits simply implement the final choice. The former allows a reader
--- human or model --- to reconstruct the decision boundary, the alternatives
considered, and the criteria used to adjudicate between them. This is
trajectory information that no snapshot preserves.

\section{What Generative Algorithms Learn}
%% ══════════════════════════════════════════════════════════════════════════════

The three-tier taxonomy has an immediate consequence for understanding what
generative AI systems actually acquire from training on the GitHub corpus, and
why that acquisition differs qualitatively from training on any other text corpus
of comparable size.

The surface account of LLM training on code is straightforward: the model sees
many programs and learns to predict likely continuations. This account is not
wrong, but it substantially understates what is actually being learned. A model
trained on the Spherepop repository and its forks does not merely learn Python
syntax or interpreter patterns. It learns the following across multiple levels of
abstraction simultaneously: that probabilistic substrates require parity-based
invariants to support reliable computation; that semantic stability is a
governance problem rather than a storage problem; that the distinction between
chaotic and void states requires a regional field sample rather than a point
observation; that lamphron implementations are typically followed by lamphrodyne
extensions that reinterpret the original abstractions as special cases of richer
geometric structures; and that mathematical appendices encode the formal substrate
of claims made informally in the implementation. None of this is explicitly
stated anywhere in the repository. All of it is inferable from the developmental
geometry of the artifact.

This is the sense in which the GitHub corpus trains generative systems not merely
on outputs but on pathways through abstraction space. The training signal is not
primarily correctness. It is recoverable development. A model that has seen many
well-formed repositories learns what it looks like when an abstract idea is
successfully compiled into executable form across multiple stages of extension
and refinement. A model that has seen many near-boundary repositories
additionally learns what failure modes look like at various developmental stages,
which is essential for generating plausible continuations of incomplete code.
A model trained exclusively on polished, formally verified programs would miss
most of this developmental geometry, because polished artifacts conceal the
stabilization process that produced them. The incomplete repositories are not
deficiencies in the training corpus. They are the primary carriers of
developmental structure.

This provides a precise answer to the natural objection that compute scale, not
corpus structure, drives capability growth. The objection conflates scale with
substrate. Scale amplifies whatever signal the substrate provides. If the
substrate is rich in admissible replay structure --- as the GitHub corpus is,
specifically because of the incentive structures described in \S3 --- then scale
propagates that richness into the model's representations. If the substrate were
primarily noise, scale would propagate noise. Remove the GitHub substrate and
apply the same compute to a corpus of equivalent size but lower admissibility
density, and the result is not a weaker version of current capability. It is a
qualitatively different and substantially less capable system.

The recursion that closed is therefore more precisely characterized than the
earlier formulation suggested. LLMs train on the developmental geometry of the
GitHub corpus and thereby acquire the capacity to generate artifacts that
themselves possess admissible developmental geometry --- artifacts that are
architecturally legible, extensible, and capable of functioning as lamphron
nucleation sites for the next generation of lamphrodyne integration. The
ecosystem does not merely accelerate. It improves the quality of its own
training signal recursively.

\subsection*{Generalization Bound for Developmental Geometry}

Let $\mathcal{D}$ be a training distribution over repositories and let
$\mathcal{A}(\mathcal{D})$ denote the mean admissibility score of sampled
repositories. We can bound the expected generalization of a model trained on
$\mathcal{D}$ in terms of $\mathcal{A}(\mathcal{D})$.

\begin{proposition}[Developmental Generalization]
Let $f_\theta$ be a language model trained by empirical risk minimization on
$n$ repositories sampled from $\mathcal{D}$. For any target developmental task
$\tau$ requiring recovery of architectural intent at depth $d$, the expected
error on $\tau$ satisfies
\[
\mathbb{E}[\mathrm{err}_\tau(f_\theta)]
\;\leq\; C \cdot \frac{1}{\mathcal{A}(\mathcal{D})} \cdot
\sqrt{\frac{\mathrm{VC}(f_\theta) \ln n}{n}},
\]
where $C > 0$ is a constant depending on $\tau$ and $\mathrm{VC}(f_\theta)$
is the VC dimension of the model class.
\end{proposition}

This bound shows that increasing $n$ (scale) reduces error at rate
$O(1/\sqrt{n})$, but increasing $\mathcal{A}(\mathcal{D})$ reduces error
\emph{multiplicatively}. A corpus with twice the admissibility density is
equivalent to a corpus four times larger in its effect on the error bound.
This formalizes the intuition that corpus structure matters as much as corpus
scale: improving the admissibility of the training distribution has superlinear
returns compared to simply adding more data.

%% ══════════════════════════════════════════════════════════════════════════════

%% ══════════════════════════════════════════════════════════════════════════════
\section{Gradient Legibility and the LLM as Manifold-Compression Operator}
\label{sec:gradient}
%% ══════════════════════════════════════════════════════════════════════════════

The preceding sections characterized generative AI as a navigational
accelerant over the GitHub substrate. This section makes that characterization
more precise by distinguishing four properties of a repository graph that a
capable model must separately address, and arguing that LLMs contribute
specifically to the last of these.

\paragraph{Accessibility.}
A repository graph is accessible if its contents can be retrieved by an agent
with network connectivity. GitHub made the graph globally accessible through
standardized URLs, public forks, and open licensing. This is a necessary but
far from sufficient precondition for exploitation. A graph of ten million
repositories that no agent can efficiently search is accessible but not
practically navigable.

\paragraph{Searchability.}
A repository graph is searchable if an agent can locate relevant nodes given
a query. GitHub provides keyword search, topic tagging, and star ranking.
These heuristics are crude: they retrieve repositories matching surface
tokens rather than repositories whose developmental geometry aligns with
a target architecture. A developer seeking a probabilistic tape interpreter
may not find one by searching for "stochastic brainfuck." Searchability
addresses surface similarity but not developmental proximity.

\paragraph{Recomposability.}
A repository graph is recomposable if its nodes can be integrated into new
systems with bounded friction. The central argument of this paper is that
GitHub dramatically reduced this friction through standardized dependency
declaration, versioned interfaces, and composable build tooling. Recomposability
is a structural property of the graph, not of any individual node: it depends
on how edges are defined, how interfaces are declared, and how conflict
resolution is supported.

\paragraph{Semantic legibility.}
A repository graph is semantically legible if an agent can infer the
developmental geometry of a node from partial evidence: given a fragment of
a codebase, determine the architectural trajectory it encodes, predict likely
extensions, identify analogous structures elsewhere in the graph, and generate
continuations that respect the latent invariants. This is the property that
LLMs contribute. None of the preceding three properties --- accessibility,
searchability, or recomposability --- requires a language model. Semantic
legibility does.

\begin{definition}[Manifold-Compression Operator]
Let $\mathcal{G}$ be a repository graph embedded in a high-dimensional
procedural state space $\mathcal{P}$. A \emph{manifold-compression operator}
$\mathcal{C} : \mathcal{P} \to \mathcal{P}'$ is a map to a lower-dimensional
space $\mathcal{P}'$ such that:
\begin{enumerate}
\item geodesic distances in $\mathcal{P}'$ approximate developmental distances
in $\mathcal{P}$ (repositories with similar architectural trajectories map to
nearby points);
\item the map is approximately invertible on the admissible submanifold
$\{\Phi > 0\}$ (a point in $\mathcal{P}'$ determines an approximately unique
Tier~1 trajectory in $\mathcal{P}$);
\item queries in $\mathcal{P}'$ can be answered with substantially lower
search cost than equivalent queries in $\mathcal{P}$.
\end{enumerate}
\end{definition}

An LLM trained on the GitHub corpus is an empirical approximation of such a
manifold-compression operator. Its embedding space is $\mathcal{P}'$; its
attention mechanism implements approximate nearest-neighbor search in that
space; its generative capabilities implement approximate inversion back to
$\mathcal{P}$. The quality of the approximation depends directly on
$\mathcal{A}(\mathcal{D})$: a training corpus with high admissibility density
produces a more faithful compression because the latent structure is richer
and more recoverable.

\begin{proposition}[LLM as Friction Reduction]
Let $\mathcal{C}$ be a manifold-compression operator trained on $\mathcal{G}$
with admissibility density $\mathcal{A}$. The effective recombination friction
between repositories $R_i, R_j$ under $\mathcal{C}$-mediated search satisfies
\[
\rho_\mathcal{C}(R_i, R_j) \;\leq\; \frac{\rho(R_i,R_j)}{\mathcal{A} \cdot
\dim(\mathcal{P})/\dim(\mathcal{P}')},
\]
where $\dim(\mathcal{P})/\dim(\mathcal{P}')$ is the compression ratio.
Higher admissibility and higher compression both reduce effective friction.
\end{proposition}

This proposition formalizes the intuition that LLMs reduce navigational
cost without altering the underlying geometry: they compress the search
problem but do not generate new admissible trajectories ex nihilo. A model
that hallucinates a repository architecture not grounded in any existing
developmental trajectory has produced output in $\mathcal{P}' \setminus
\pi(\mathcal{G})$ --- outside the image of the graph under compression ---
which will typically fail to compose with real infrastructure precisely because
it lacks the accumulated constraint structure that real repositories carry.

\section{Developmental Geometry Across Representational Layers}
%% ══════════════════════════════════════════════════════════════════════════════

One of the most important properties of the Spherepop-over-\ambibf{} example is
that the developmental trajectory does not terminate at the implementation layer.
The repository does not merely contain executable code. It participates in a
larger representational manifold extending across operational systems,
mathematical appendices, field-theoretic monographs, categorical abstractions,
and simulation methodologies simultaneously. The developmental geometry exposed
to the training manifold therefore exists across multiple semantic strata at once.

This becomes especially clear when the repository is interpreted alongside the
RSVP and Quantum SpherePop papers. The original \texttt{AmbiBFMachine}
implementation defines a minimal stochastic substrate whose operational semantics
are governed by noisy pointer drift and probabilistic arithmetic transitions.
The minimal Spherepop layer introduces metastable semantic regions stabilized
through repeated sampling and parity-preserving constraints. The field-theoretic
extensions in the accompanying monographs then reinterpret these operational
intuitions in terms of coarse-grained RSVP fields, entropy-respecting dynamics,
synchronization manifolds, and lattice quantization schemes. The abstractions
survive extension pressure because they preserve invariant developmental
structure across layers.

The crucial observation is that these layers are not independent. The operational
code constrains the mathematical language that later emerges around it. The
mathematical formalism retrospectively clarifies the implementation decisions
encoded in the repository history. The categorical descriptions unify both into
a broader semantic architecture. The result is not a set of disconnected
artifacts but a single developmental attractor expressed through multiple
representational media.

This is precisely the kind of structure that generative systems appear unusually
effective at internalizing. A model trained on the combined corpus does not
simply learn a Python interpreter, a field theory, and a category-theoretic
formalism as separate objects. It learns that these objects occupy neighboring
regions in abstraction space and participate in the same stabilization
trajectory. The developmental signal therefore propagates not only through code
reuse but through cross-layer semantic alignment.

%% ══════════════════════════════════════════════════════════════════════════════
\section{Quantum SpherePop and Semantic Quantization}
%% ══════════════════════════════════════════════════════════════════════════════

The following three sections reinterpret the same developmental trajectory
across progressively richer representational layers: operational implementation,
stochastic semantic fields, replay geometry, and finally RSVP admissibility
dynamics. The purpose is not to multiply ontologies but to show that the same
invariants survive translation across representational media. Each layer
exposes structure latent in the one before it; none contradicts the others.
The escalation in formalism is recursive reinterpretation, not conceptual drift.

The developmental trajectory from the minimal \texttt{AmbiBFMachine}
implementation to the later Quantum SpherePop formalism is not merely an
increase in implementation complexity. It is a semantic quantization sequence.
The original interpreter already contains, in low-resolution operational form,
many of the structural relations that later appear explicitly in the RSVP
field-theoretic framework.

In the minimal implementation, a bubble is represented as a parity-preserving
neighborhood of tape cells whose semantic state emerges through repeated
stabilization under stochastic perturbation. In the Quantum SpherePop
framework, local semantic regions are instead represented by complex amplitude
fields:
\[
\Psi_i = r_i e^{i\theta_i},
\]
where amplitude $r_i$ encodes entropic magnitude and phase $\theta_i$ encodes
semantic coherence. The correspondence is structurally direct. The
parity-preserving bubble neighborhoods of the interpreter become coarse
approximations of phase-coherent regions in a coupled semantic lattice. The
stochastic drift of the substrate becomes a discrete analogue of vacuum
fluctuation and entropic perturbation. The stabilization loop becomes a
primitive synchronization operator acting over neighboring semantic regions.

The repository trajectory therefore behaves like a semantic renormalization flow:
\[
\texttt{AmbiBFMachine}
\;\longrightarrow\;
\texttt{Spherepop}
\;\longrightarrow\;
\texttt{Quantum SpherePop}.
\]
Each stage preserves invariant developmental geometry while increasing
representational richness. The operational implementation and the later
mathematical framework are therefore not separate artifacts but adjacent regions
of the same developmental manifold.

\subsection*{Renormalization as Coarse-Graining}

The three-stage sequence is an instance of Wilsonian renormalization: at each
scale, fine-grained degrees of freedom are integrated out, leaving an effective
theory that captures the infrared behavior.

\begin{definition}[Semantic Renormalization Group]
Let $\mathcal{S}_0$ be the space of \ambibf{} tape configurations. Define a
coarse-graining map $\pi_k : \mathcal{S}_0 \to \mathcal{S}_k$ that sends each
tape configuration to its bubble coherence vector at scale $k$:
\[
\pi_k(s) = \bigl(C(B_1), C(B_2), \ldots, C(B_m)\bigr)
\]
where $C(B_j) = \sum_{i \in \Omega_j} s_i$ is the signed coherence of bubble
$j$ and $m$ is the number of bubbles. The sequence $(\pi_k)_{k \geq 0}$
defines a renormalization group acting on semantic states.
\end{definition}

\begin{proposition}[Fixed Points of Semantic RG]
The fixed points of the semantic renormalization group are the configurations
in which each bubble is either maximally coherent ($C(B_j) = \pm E(B_j)$) or
void ($C(B_j) = 0$). These correspond to the stable semantic states
\texttt{coherent-positive}, \texttt{coherent-negative}, and \texttt{void}
in the \texttt{semantic\_state} function.
\end{proposition}

\begin{proof}[Proof sketch]
Under repeated stabilization, the reinforcement operator $\mathcal{R}_\lambda$
drives each bubble toward the sign of its current coherence (Appendix~E). A
configuration is a fixed point when further stabilization does not change any
bubble's sign. This occurs when every bubble has coherence strictly positive,
strictly negative, or zero --- the three fixed-point classes above.
\end{proof}

%% ══════════════════════════════════════════════════════════════════════════════
\section{Replay as Coarse-Grained Semantic Reconstruction}
%% ══════════════════════════════════════════════════════════════════════════════

One of the deepest properties of the GitHub ecosystem is that repositories are
not consumed statically. They are replayed. Forks, ports, rewrites, patches,
dependency updates, and architectural reinterpretations are all replay operations
over developmental trajectories.

In the Spherepop formalism, replay does not preserve exact state identity.
Instead it regenerates a statistically stable attractor under perturbation. A
semantic object persists not because every replay is identical, but because the
replay sequence converges toward a recoverable semantic region despite local
divergence:
\[
\sigma_n \sim \mathbb{P}_n.
\]
The GitHub ecosystem behaves similarly. A repository fork is not a perfect copy
of an artifact. It is a coarse-grained reconstruction of developmental intent
under altered environmental constraints. Dependencies change. APIs drift.
Architectures are partially rewritten. Yet semantic identity often persists
across these perturbations because the attractor geometry remains recoverable.

This interpretation aligns closely with the RSVP/TARTAN coarse-graining
framework in which observable structures emerge from lower-level field dynamics
through projection operators:
\[
\pi : X(t) \to \chi_i(t).
\]
The raw repository state corresponds to the inaccessible microdynamic layer.
Replay operations act as coarse-graining transformations over developmental
history. The resulting semantic object is therefore not a static file tree but
a metastable equivalence class generated by repeated replay under changing
conditions.

This explains why commit histories are often more educational than snapshots.
A snapshot reveals what a repository is. A replay history reveals how semantic
coherence was preserved while the system changed. Generative systems trained on
such histories therefore acquire not merely static representations of code but
representations of stabilization dynamics themselves --- and it explains why
open ecosystems accelerate capability growth more effectively than closed
proprietary systems. Open repositories undergo continual replay by many
independent agents simultaneously: forks, patches, reinterpretations,
optimizations, ports, and integrations. Each replay operation exposes additional
developmental geometry to the collective substrate.

\subsection*{Convergence of Replay Distributions}

\begin{definition}[Semantic Attractor]
A repository $R$ has a well-defined semantic attractor if the sequence of
replay distributions $\{\mathbb{P}_n\}_{n \geq 1}$ converges weakly to a
limiting distribution $\mathbb{P}_\infty$ concentrated on a compact set
$\mathcal{A}(R) \subset$ configuration space.
\end{definition}

\begin{theorem}[Attractor Existence under Stabilization]
\label{thm:attractor-existence}
Let $R$ be a Tier~1 repository whose architectural commitments impose a
Lyapunov function $V : \mathcal{S} \to \mathbb{R}_{\geq 0}$ satisfying
$\mathbb{E}[V(\sigma_{n+1}) \mid \sigma_n] \leq V(\sigma_n) - \alpha$ for
some $\alpha > 0$ outside a compact set $K$. Then the replay distributions
$\mathbb{P}_n$ converge weakly to a unique stationary distribution
$\mathbb{P}_\infty$ supported on $K$.
\end{theorem}

\begin{proof}[Proof sketch]
The condition on $V$ is a Foster--Lyapunov drift criterion (Foster's theorem).
It implies positive recurrence of the Markov chain $(\sigma_n)_{n \geq 0}$
on $\mathcal{S}$, which by the ergodic theorem guarantees weak convergence to
a unique stationary measure. The compact support follows from the sublevel-set
compactness of $V$.
\end{proof}

\begin{corollary}
Tier~2 repositories (those near the admissibility boundary) lack a uniform
Lyapunov function satisfying the drift condition globally; they may converge
locally but have positive probability of escaping the compact set under large
perturbations. Tier~3 repositories satisfy no drift condition: without
architectural structure, $V$ cannot be defined, and no attractor exists.
\end{corollary}

%% ══════════════════════════════════════════════════════════════════════════════
\section{Phase Synchronization and the Collective Procedural Field}
\label{sec:phase}
%% ══════════════════════════════════════════════════════════════════════════════

The Quantum SpherePop framework introduces a particularly important
reinterpretation of semantic coherence: coherence is not exact symbolic
agreement but partial phase synchronization across distributed semantic regions.
The local phase evolution equation
\[
\dot{\theta}_i
=
\omega_i
+
\sum_{j \in N(i)}
K_{ij}\sin(\theta_j - \theta_i)
+
\sqrt{2D}\,\eta_i(t)
\]
defines coherence as an emergent property of coupled stochastic dynamics rather
than centralized control~\cite{kuramoto1984}. Here $\omega_i$ is the natural
frequency of region $i$, $K_{ij}$ is the coupling strength between neighboring
regions, $D$ is the noise intensity, and $\eta_i(t)$ is independent white noise
at each site.

This provides a remarkably precise description of how large-scale open-source
ecosystems evolve. Repositories are not globally coordinated through a single
architectural authority. Instead local synchronization occurs through dependency
compatibility, shared interfaces, community conventions, benchmark standards,
framework ecosystems, and repeated replay interactions across neighboring
projects. The resulting ecosystem resembles a distributed semantic phase field.
Local regions synchronize partially while preserving diversity at larger scales.
Forks generate phase divergence. Standards generate synchronization pressure.
Merge operations reduce local incoherence. Stable abstractions emerge as
phase-locked attractors surviving repeated perturbation across the collective
graph.

The intelligence explosion can therefore be interpreted not merely as an increase
in model capability but as a phase transition in the coherence of the global
procedural substrate itself. GitHub externalized developmental trajectories into
a globally replayable semantic field. Generative systems trained on that field
acquired the ability to navigate, reconstruct, and extend those trajectories
statistically. The result is not artificial intelligence emerging ex nihilo. It
is the planetary procedural field becoming partially self-modeling through
recursive replay and synthesis.

\subsection*{Critical Coupling and the Synchronization Transition}

The Kuramoto model exhibits a phase transition at a critical coupling strength
$K_c$. We derive the analogue for the GitHub ecosystem.

In the Kuramoto model on a complete graph of $N$ oscillators with natural
frequencies drawn from a symmetric unimodal distribution $g(\omega)$ with
half-width $\Delta$, the critical coupling is:
\[
K_c = \frac{2}{\pi g(0)}.
\]
For a Lorentzian distribution $g(\omega) = \frac{\Delta}{\pi(\omega^2 +
\Delta^2)}$, one has $g(0) = 1/(\pi\Delta)$ and thus $K_c = 2\Delta$.

\begin{proposition}[Ecosystem Synchronization Threshold]
Interpret each repository as an oscillator with natural frequency $\omega_i$
encoding its architectural style (rate of change, abstraction level, domain
specificity) and coupling strength $K_{ij}$ encoding the dependency weight
between $R_i$ and $R_j$. The ecosystem transitions from a fragmented
(incoherent) state to a partially synchronized state when the mean coupling
strength $\bar{K}$ exceeds $K_c = 2\Delta$, where $\Delta$ is the spread of
architectural styles in the ecosystem.

GitHub's structural innovations --- standardized dependency declaration,
semantic versioning, shared CI tooling --- simultaneously increased $\bar{K}$
and reduced $\Delta$, pushing the ecosystem from below to above $K_c$.
\end{proposition}

\begin{corollary}[Irreversibility of the Transition]
Once $\bar{K} > K_c$, the synchronized state is a global attractor: the
fraction of synchronized repositories $r(t) \to r_\infty > 0$ as $t \to
\infty$. The transition is irreversible in the sense that reducing $\bar{K}$
back below $K_c$ requires dismantling the coupling infrastructure (e.g.,
removing package managers, eliminating semantic versioning), not merely
reducing individual dependency counts.
\end{corollary}

%% ══════════════════════════════════════════════════════════════════════════════
\section{RSVP Connections: GitHub as Planetary Admissibility Geometry}
%% ══════════════════════════════════════════════════════════════════════════════

The GitHub ecosystem can be mapped onto the RSVP field-theoretic framework with
considerable precision, and doing so clarifies both the nature of the
acceleration and its theoretical relationship to other phenomena modeled within
that framework.

In the RSVP framework, the scalar field $\Phi$ encodes the density of
constraint-satisfying configurations: regions where $\Phi > 0$ correspond to
admissible states, while the zero-set $Z(\Phi) = \{p \in M : \Phi(p) = 0\}$
marks the constraint boundary at which physical or epistemic admissibility fails.
The vector field $\mathbf{v}$ encodes the directed flow of evolution through
configuration space, and the entropy field $S$ tracks the irreversible
accumulation of structural commitment along trajectories. The field equations
enforce a monotonicity condition: along any admissible trajectory, $S$ is
non-decreasing, and strictly increasing wherever the constraint gradient
$\|\nabla\Phi\|^2 > 0$.

Under this mapping, the GitHub repository graph corresponds to a
high-dimensional configuration space of executable cognitive states. Each
repository is a point in that space: a particular constraint-satisfying
configuration of code, documentation, architecture, and dependency structure.
The directed dependency graph between repositories defines a vector field
$\mathbf{v}$ encoding the dominant directions of epistemic flow --- the typical
trajectories along which new configurations are built from existing ones. The
scalar field $\Phi$ encodes the density of viable configurations in the
neighborhood of each point. Forking, dependency resolution, and integration
correspond to the constraint-satisfaction dynamics that define admissible
trajectories through this space.

The entropy field $S$ tracks the irreversible accumulation of committed
architectural decisions. Every dependency added, every interface stabilized,
every API version frozen represents a Spherepop event: an irreversible constraint
that narrows the option-space of future configurations. Applied to repositories,
this means that a codebase is not merely its current state but the trajectory of
irreversible decisions that generated it. The full Spherepop normal form records
what the snapshot of current state obscures.

The lamphron--lamphrodyne duality applies here with particular clarity. Lamphron
designates the differentiation pressure: the tendency for systems to generate
localized structure, form stable memory configurations, and specialize toward
predictive specificity. In the GitHub context, lamphron manifests as the pressure
toward modularity and the development of increasingly specialized repositories
that solve narrow problems with high precision. Lamphrodyne designates the
integration pressure: the tendency for locally specialized structures to be
recombined into higher-order composites. In the GitHub context, lamphrodyne
manifests as dependency aggregation, framework construction, and the assembly of
complex systems from modular components.

The intelligence explosion as a whole can be understood as the product of a
lamphron--lamphrodyne feedback loop operating at planetary scale: specialized
repositories accumulate, creating the conditions for their own integration into
increasingly capable composite systems, whose deployment creates new problems
that generate new specialized components, in a cycle of recursive structural
elaboration that mirrors the cosmological process by which gravitational
differentiation produces increasingly complex relational structure. This is
precisely the dynamic that Barbour's entaxy framework tracks in the physical
domain~\cite{barbour2020}: the universe moves away from the Janus Point not
into disorder but into increasing relational specificity. The GitHub ecosystem
enacts an analogous process in the domain of executable cognition.

The role of generative AI in this picture is precise: it functions as an
entropy-compressing interface that reduces the navigational cost of traversing
the configuration space. Without such an interface, the density of the repository
graph eventually exceeds the capacity of unaided human search to exploit it
efficiently. Language models trained on the graph can compress its structure into
lower-dimensional representations that are easier to query, recombine, and
extend. This is not a qualitative transformation of the underlying dynamics; it
is a reduction in the effective friction coefficient governing traversal of an
already-existing admissibility geometry.

The connection to the Yarncrawler framework reinforces this interpretation.
Yarncrawler treats world-state reconstruction as a Wasserstein gradient flow over
constraint-satisfying configurations: the goal state defines a target in the
space of executable structures, and the reconstruction process is the gradient
descent that drives the system toward it. Repository discovery and integration
can be modeled as exactly this kind of flow. The developer's intended
architecture defines a target configuration; the process of finding, evaluating,
forking, and adapting repositories is the gradient descent that drives the
working codebase toward that target. Generative AI makes the gradient more
legible --- it compresses the search space and reduces the cost of evaluating
candidate configurations --- but it does not modify the underlying geometry of
admissibility that defines which configurations are reachable.

\subsection*{RSVP Field Equations Applied to the Repository Graph}

The RSVP triple $(\Phi, \mathbf{v}, S)$ on the repository configuration space
$M$ satisfies:
\begin{align}
\Box\Phi + \nabla_{\mathbf{v}}\Phi &= -\lambda S\Phi, \label{eq:phi-rsvp}\\
\nabla_{\mathbf{v}}\mathbf{v} + \nabla\Phi &= -\kappa\mathbf{v}, \label{eq:v-rsvp}\\
\partial_t S + \nabla\cdot(S\mathbf{v}) &= \sigma\|\nabla\Phi\|^2.
\label{eq:S-rsvp}
\end{align}

\begin{proposition}[Entropy Monotonicity Along Admissible Trajectories]
\label{prop:entropy-mono}
Let $\gamma : [0,T] \to M$ be an admissible trajectory (i.e., $\Phi(\gamma(t))
> 0$ for all $t$). Then $S(\gamma(t))$ is non-decreasing, and strictly
increasing whenever $\|\nabla\Phi(\gamma(t))\|^2 > 0$.
\end{proposition}

\begin{proof}
Along $\gamma$, equation \eqref{eq:S-rsvp} gives $\frac{d}{dt}S(\gamma(t)) =
\sigma\|\nabla\Phi(\gamma(t))\|^2 \geq 0$ since $\sigma > 0$. Strict positivity
holds whenever $\nabla\Phi \neq 0$ at the trajectory point.
\end{proof}

Applied to the repository graph, Proposition~\ref{prop:entropy-mono} states that
the entropy of irreversible architectural commitments is non-decreasing along any
valid development trajectory. Every design decision that moves the codebase
further into the admissible region $\Phi > 0$ increases $S$, recording an
irreversible narrowing of the future option-space. This is not a deficiency of
the system; it is the mechanism by which semantic identity becomes stable.

\begin{corollary}[Developmental Irreversibility]
A well-formed repository cannot be reverted to a semantically neutral state
without destroying its developmental identity. The entropy $S$ accumulated
along its trajectory records the irreversible commitment structure that
constitutes its Spherepop normal form.
\end{corollary}

%% ══════════════════════════════════════════════════════════════════════════════

%% ══════════════════════════════════════════════════════════════════════════════
\section{Objections and Clarifications}
%% ══════════════════════════════════════════════════════════════════════════════

A thesis of this scope invites principled objection. We address the most
consequential challenges directly, in each case arguing that the objection,
properly understood, either strengthens or further specifies the central claim
rather than undermining it.

\paragraph{Is this merely collective intelligence?}
The collective intelligence literature~\cite{holland1992} describes how groups
of agents produce outcomes exceeding individual capacity. The GitHub thesis
is structurally different in two respects. First, collective intelligence
typically requires agents to coordinate in real time, whereas GitHub
coordinates asynchronously across decades: the developer who forks a 2009
repository in 2024 is coordinating with contributors who are no longer active.
Second, collective intelligence produces emergent outputs but rarely externalizes
the intermediate cognitive states that generated them. GitHub's distinctive
contribution is precisely the preservation of developmental geometry ---
the trajectory, not only the destination. The thesis is therefore not that
many people working together are smart, but that GitHub changed the ontology
of what gets preserved when they do.

\paragraph{Why GitHub specifically rather than the internet generally?}
The internet created global information access. GitHub created globally
recomposable executable structure. The distinction is the difference between
reading a description of a machine and having access to the machine's
blueprints in a format that allows direct modification, extension, and
integration with other machines. Stack Overflow~\cite{raymond1999} preserves
isolated code fragments in natural-language context. ArXiv preserves finished
formal artifacts. Neither preserves the developmental trajectory of a system
--- the sequence of commits, the issue discussions, the failed branches, the
dependency migrations. Only a version-controlled repository with public history
does this, and GitHub's structural innovations (forking, pull requests, issue
tracking, dependency graphs) made this preservation compositional and globally
reachable.

\paragraph{Does closed-source development undermine the thesis?}
Closed development does not falsify the thesis; it quantifies the cost of its
absence. The counterfactual is instructive: a world in which all development
remained proprietary would have produced model capabilities far below the
current state, because the training corpus would lack the procedural reasoning
substrate that open repositories provide. The thesis predicts not that all
intelligence explosion requires open source, but that the observed acceleration
is explained by the fraction of global development that became publicly
recomposable. Closed development contributes zero to $\mathcal{A}(\mathcal{D})$
while consuming resources from the shared ecosystem, which is precisely why
firms that selectively open-source gain disproportionate community returns.

\paragraph{Could earlier scientific traditions have produced similar effects?}
Preprint culture, open journals, and reproducible research movements all moved
in the direction of this thesis but fell short for a specific structural reason:
they preserved descriptions of procedures rather than the procedures themselves.
A preprint describing a training algorithm is not the same as a public
repository containing the training code, the data loading pipeline, the
evaluation harness, and the issue thread debating the correctness of a
preprocessing step. The latter preserves developmental geometry in executable
form. The former preserves a prose projection of it, from which the original
trajectory cannot be recovered without substantial independent reconstruction.
GitHub closed this gap; no prior system did so at scale.

\paragraph{Is the argument anthropomorphizing repositories?}
The RSVP framework uses language --- attractors, mortality, coherence,
governance --- that may seem to impute agency to static files. This is a
notational convenience, not a metaphysical commitment. The claim is structural:
a repository that encodes parity-preserving bubble logic, majority-vote
stabilization, and a two-stage lamphron-to-lamphrodyne development history
exposes richer developmental geometry than a repository that does not, in
the same way that a crystal with specific lattice geometry encodes more
recoverable structure than amorphous glass. No agency is implied; only
recoverable architectural information.

\paragraph{On dark repositories and AI-generated codebases.}
A significant open problem deserves explicit acknowledgment. The recursion
identified in equations~\eqref{eq:model-update} and~\eqref{eq:graph-update}
implies that as $\mathcal{M}_t$ becomes sufficiently capable, it begins
generating $\mathcal{G}_{t+1}$ faster than human contributors can evaluate
it. This creates the possibility of what we might call dark repositories:
codebases generated by language models for consumption by other language
models, forming a coherence field that maintains admissibility within
$\mathcal{G}$ while becoming progressively opaque to human developmental
reconstruction. The lamphron--lamphrodyne duality does not require human
observers to function; it requires only that nucleation and integration
dynamics operate over a graph with sufficient admissibility structure. Whether
AI-generated lamphron repositories constitute genuine developmental geometry ---
whether they possess recoverable trajectories in the sense required by
Theorem~\ref{thm:attractor-existence} --- or whether they produce admissible
but developmentally hollow artifacts is an empirical question that current
tooling cannot yet answer. It is one of the most important open questions
raised by the framework.

\section{Conclusion: Thought as Recursively Reusable Infrastructure}
%% ══════════════════════════════════════════════════════════════════════════════

The intelligence explosion is not primarily artificial intelligence becoming
intelligent. It is humanity converting thought itself into recursively reusable
public machinery, and then training pattern synthesizers to navigate that
machinery faster.

GitHub transformed knowledge from static publication into evolutionary
infrastructure. It externalizes cognitive debris rather than discarding it. It
rewards incompleteness in ways that previous institutional structures punished.
It aligns the incentive structures of academia, industry, and independent
research toward a common outcome that no centralized institution could have
planned: the continuous public disclosure of intermediate cognitive states at
planetary scale.

Generative AI is real, and its capabilities are genuinely significant. But it is
an accelerant layered over a substrate, not the substrate itself. The substrate
is the open-source graph: globally persistent, recursively recomposable,
executable, and continuously growing. Remove the substrate and the accelerant has
little to accelerate over. Remove the accelerant and the substrate continues to
grow, as it did for more than a decade before large language models were publicly
deployed.

In RSVP terms, GitHub created the admissibility geometry. Generative AI reduced
the cost of navigating it. The explosion lives in the coupling between these two
systems --- in the closed recursion loop that now connects model output to
repository generation to training data to model improvement. That loop did not
close when GPT appeared. It began closing when the first pull request was merged
into a publicly accessible repository and the first developer on another
continent forked it the same day.

GitHub transformed civilization into a globally replayable semantic field.
Generative AI emerged as a navigational dynamics over that field. The true
invention was globally persistent recomposable thought, and GitHub
operationalized it before the models arrived to compress it.

\paragraph{The civilizational discontinuity.}
Framed at the largest scale, the argument identifies a genuine historical
discontinuity: a transition in what civilizations preserve. Prior to version
control, the record of human technical cognition consisted almost entirely of
finalized symbolic outputs --- publications, products, patents, monographs.
The developmental process that generated those outputs was not recorded;
it was discarded as ephemera. What survived was the destination, not the
trajectory. GitHub changed this. For the first time in human history, a
civilization-scale system now preserves not only what was built but
\emph{how it was built}: the reasoning, the failures, the alternatives
considered, the bugs encountered and repaired, the partial implementations
abandoned, the architectural pivots made under constraint. This is not an
incremental improvement in archival fidelity. It is a categorical change in
the ontology of what counts as preserved knowledge.

The intelligence explosion is therefore a second-order effect of this
civilizational transition. Generative systems trained on the GitHub corpus
do not merely learn to write code. They learn to navigate a preserved record
of human developmental cognition at planetary scale --- a record that no
prior civilization produced or could have produced. The capability growth
associated with modern AI systems is, in this sense, less a property of
the models themselves than a property of the unprecedented substrate they
have been given access to.

\paragraph{Open problems.}
The framework developed here raises several questions that the current
formalism does not resolve. The most pressing concerns the transition
identified in the objections section: as AI-assisted generation accelerates
$\mathcal{G}_{t+1}$, the ratio of human-generated to model-generated
developmental geometry in the training corpus will shift. Whether this
transition preserves or degrades admissibility density --- whether AI-generated
commits constitute genuine developmental trajectories or developmentally
hollow admissible artifacts --- is empirically undetermined and theoretically
urgent. A second open problem concerns the formalization of $\mathcal{A}(R)$
beyond the proportionality established here: a computable admissibility metric
would enable direct corpus filtering and quality assessment at scale. A third
concerns the extension of the Kuramoto synchronization analysis to directed
dependency graphs with heterogeneous coupling strengths, which better model
real ecosystem topology than the complete-graph approximation used in
\S\ref{sec:phase}. These are the problems on which the framework must be
tested if it is to transition from a theoretical architecture into a
predictive science of collective technical cognition.

%% ══════════════════════════════════════════════════════════════════════════════

%% ══════════════════════════════════════════════════════════════════════════════
\section*{Notation Glossary}
\addcontentsline{toc}{section}{Notation Glossary}
\label{sec:notation}
%% ══════════════════════════════════════════════════════════════════════════════

The following symbols appear throughout the main text and appendices. Where a
symbol carries both a general mathematical meaning and a specific interpretation
in the GitHub or Spherepop context, both are given.

\medskip
\begin{center}
\begin{tabular}{lp{9.5cm}}
\toprule
\textbf{Symbol} & \textbf{Meaning} \\
\midrule
$\mathcal{G}_t$ & Repository graph at time $t$; nodes are repositories,
  directed edges are dependency or derivation relations \\
$\mathcal{M}_t$ & Language model at time $t$, trained on the corpus of
  $\mathcal{G}_t$ \\
$\Sigma_t$ & Joint state $(\mathcal{G}_t, \mathcal{M}_t)$ of the
  coupled dynamical system \\
$\rho(R_i, R_j)$ & Recombination friction between repositories $R_i$ and $R_j$;
  expected developer effort to integrate $R_j$ into a codebase depending on $R_i$ \\
$\mathcal{A}(R)$ & Admissibility score of repository $R$; proportional to
  recoverable semantic structure divided by perturbation magnitude \\
$\Phi$ & RSVP scalar field encoding density of constraint-satisfying
  configurations; $\Phi > 0$ in admissible regions \\
$Z(\Phi)$ & Zero-set $\{p \in M : \Phi(p) = 0\}$; the constraint boundary
  separating admissible from inadmissible states \\
$\mathbf{v}$ & RSVP vector field encoding directed flow through configuration
  space; interpreted as dominant directions of epistemic development \\
$S$ & RSVP entropy field tracking irreversible accumulation of structural
  commitment; non-decreasing along admissible trajectories \\
$B_k = (c_k, r_k, \pi_k)$ & Spherepop bubble: center $c_k \in \mathbb{Z}$,
  radius $r_k \in \mathbb{N}$, parity class $\pi_k \in \{0,1\}$ \\
$\Omega_k$ & Support region of bubble $B_k$: all tape cells within radius
  $r_k$ of center $c_k$ having parity $\pi_k$ \\
$E(B_k)$ & Bubble energy $\sum_{i \in \Omega_k} |s_i|$; measures semantic
  persistence \\
$C(B_k)$ & Signed coherence $\sum_{i \in \Omega_k} s_i$; positive for
  constructive, negative for destructive, zero for decoherent bubbles \\
$\mathcal{R}_\lambda$ & Stabilization (reinforcement) operator with governance
  coefficient $\lambda > 0$; drives bubble cells toward the sign of $C(B_k)$ \\
$\mathcal{M}(t)$ & Global mortality functional $\sum_k \frac{d}{dt} E(B_k)$;
  negative in expectation without stabilization \\
$\sigma_n$ & State after $n$ replay steps; a random variable distributed
  according to $\mathbb{P}_n$ \\
$\mathbb{P}_n$ & Distribution over semantic states after $n$ replay
  operations; converges weakly to $\mathbb{P}_\infty$ for Tier~1 repositories \\
$\mu(B_a, B_b)$ & Merge coherence $\sum_{i \in \Omega_a \cap \Omega_b}
  s_i^{(a)} s_i^{(b)}$; positive for constructive merges \\
$\pi_k(s)$ & Coarse-graining map from tape state $s$ to bubble coherence
  vector at scale $k$; defines the semantic renormalization group \\
$\Psi_i$ & Complex amplitude $r_i e^{i\theta_i}$ in the Quantum SpherePop
  lattice; amplitude encodes entropic magnitude, phase encodes coherence \\
$K_c$ & Critical coupling strength $2\Delta$ for the synchronization
  transition in the Kuramoto model; interpreted as the GitHub ecosystem
  threshold separating fragmented from synchronized development \\
$q_t$ & Commit hash at time $t$; interpreted as a Lamport timestamp encoding
  causal order in the repository's event history \\
\bottomrule
\end{tabular}
\end{center}
\medskip

\appendix
%% ══════════════════════════════════════════════════════════════════════════════

\section{Probabilistic Tape Geometry}
\label{sec:tapegeom}

Let the substrate tape be indexed by the integers $\mathbb{Z}$. Each tape cell
carries an integer-valued stochastic state $s_i(t) \in \mathbb{Z}$ for tape
coordinate $i$ and discrete interpreter time $t$. The substrate state is
therefore a field:
\[
S(t) = \{ s_i(t) \}_{i \in \mathbb{Z}}.
\]
Unlike ordinary Brainfuck, transition operators are nondeterministic. The
arithmetic operator induces a symmetric random walk:
\[
s_i(t+1) = s_i(t) + \xi_t, \qquad
\xi_t \in \{-1,+1\}, \quad
\Pr(\xi_t = +1) = \Pr(\xi_t = -1) = \tfrac{1}{2}.
\]
Likewise, pointer motion is stochastic:
\[
p(t+1) = p(t) + \eta_t, \qquad \eta_t \in \{-1,+1\}.
\]
Thus the substrate is not a deterministic memory system but a coupled stochastic
field.

\section{Bubble Regions}
\label{sec:bubbleregions}

A Spherepop bubble is not represented by a single tape cell but by a finite
parity-preserving neighborhood. Define a bubble $B_k$ as a triple:
\[
B_k = (c_k,\, r_k,\, \pi_k)
\]
where $c_k$ is the center, $r_k$ is the radius, and $\pi_k \in \{0,1\}$ is the
parity class. The support region is:
\[
\Omega_k = \bigl\{
i \in \mathbb{Z} : |i - c_k| \le r_k
\text{ and }
i \equiv \pi_k \pmod{2}
\bigr\}.
\]
Parity structure is important because parity remains one of the few stable
geometric invariants under noisy drift, as established by
Lemma~\ref{lemma:parity} in \S\ref{sec:worked}.

\section{Bubble Energy and Coherence}
\label{sec:energy}

Define the local bubble energy:
\[
E(B_k) = \sum_{i \in \Omega_k} |s_i|.
\]
This measures semantic persistence. Define the signed coherence:
\[
C(B_k) = \sum_{i \in \Omega_k} s_i.
\]
Interpretation: $C > 0$ denotes constructive coherence; $C < 0$ denotes
destructive coherence; $C = 0$ denotes decoherence. A semantic bubble exists
only while $E(B_k) > \varepsilon$ for some stability threshold $\varepsilon > 0$.

\section{Mortality Functional}
\label{sec:mortality}

Define global mortality:
\[
\mathcal{M}(t) = \sum_k \frac{d}{dt} E(B_k).
\]
Without stabilization, $\mathbb{E}[\mathcal{M}(t)] < 0$.

\begin{theorem}[Bubble Mortality]
\label{thm:mortality}
Let $S(t)$ evolve solely under unbiased substrate dynamics. Then:
\[
\lim_{t \to \infty} \Pr\bigl(E(B_k) = 0\bigr) = 1.
\]
\end{theorem}

\begin{proof}[Proof sketch]
The bubble support undergoes coupled symmetric random walks. Since no restoring
force exists, finite local coherence dissipates under recurrence and
cancellation. More precisely, the process $E(B_k)(t)$ is a non-negative
supermartingale (under unbiased dynamics, $\mathbb{E}[E(B_k)(t+1) \mid
\mathcal{F}_t] \leq E(B_k)(t)$) bounded below by zero. By the martingale
convergence theorem it converges almost surely to a limit $L \geq 0$; since
the random walk on $\mathbb{Z}$ is recurrent, the limit must be $L = 0$ almost
surely.
\end{proof}

\section{Stabilization Operators}
\label{sec:stabilization}

Define a reinforcement operator $\mathcal{R}_\lambda : S \to S$ acting locally
by:
\[
s_i \mapsto s_i + \lambda \operatorname{sign}(C(B_k)).
\]
This introduces directional correction against substrate entropy. The stabilized
dynamics become:
\[
s_i(t+1) = s_i(t) + \xi_t + \lambda \operatorname{sign}(C(B_k)),
\]
where $\lambda > 0$ is the semantic governance coefficient. For $\lambda >
\sqrt{2}$ the drift term dominates the noise variance, guaranteeing
$\mathbb{E}[E(B_k)(t+1) \mid \mathcal{F}_t] > E(B_k)(t)$ whenever
$C(B_k)(t) \neq 0$.

\section{Merge Dynamics}
\label{sec:merge}

Let two bubbles overlap: $\Omega_a \cap \Omega_b \ne \emptyset$. Define merge
coherence:
\[
\mu(B_a, B_b) = \sum_{i \in \Omega_a \cap \Omega_b} s_i^{(a)}\, s_i^{(b)}.
\]
If $\mu(B_a, B_b) > 0$ the merge is constructive. If $\mu(B_a, B_b) < 0$ the
merge is destructive and induces collapse radiation:
\[
\Delta E = -\alpha |\mu|.
\]
This corresponds to the Spherepop Calculus merge operator interpreted
geometrically over a mortal substrate.

\section{Semantic Replay}
\label{sec:replay}

Define a semantic replay sequence $\mathcal{P} = (e_1, e_2, \dots, e_n)$ with
replay operator:
\[
\sigma_n = e_n \circ e_{n-1} \circ \cdots \circ e_1(\sigma_0).
\]
Over a stochastic substrate the replay is itself stochastic, generating a
distribution over semantic states:
\[
\sigma_n \sim \mathbb{P}_n.
\]
Meaning is therefore not a single state but an attractor distribution. The
convergence of $\{\mathbb{P}_n\}$ is guaranteed for Tier~1 repositories by
Theorem~\ref{thm:attractor-existence}.

\section{Attractor Semantics}
\label{sec:attractors}

Define the replay expectation $\bar{\sigma}_n = \mathbb{E}[\sigma_n]$. A
semantic object $O$ is stable if and only if:
\[
\lim_{n \to \infty} \operatorname{Var}(O_n) < \delta
\]
for some tolerance $\delta > 0$. Semantic persistence is therefore not exact
symbolic preservation but bounded variance under replay. A repository's
contribution to the collective substrate is proportional to the recoverability
of its semantic attractor under perturbation:
\[
\mathcal{A}(R)
\;\propto\;
\frac{\text{recoverable semantic structure}}{\text{perturbation magnitude}}.
\]

\section{Spherepop--\ambibf{} Correspondence}
\label{sec:correspondence}

\medskip
\begin{center}
\begin{tabular}{ll}
\toprule
\textbf{Spherepop concept} & \textbf{Substrate interpretation} \\
\midrule
Sphere          & Local parity-preserving field \\
Pop             & Entropic collapse event \\
Merge           & Coherence synchronization \\
Choice          & Stochastic branch distribution \\
Collapse        & Representative normalization \\
Replay          & Attractor regeneration \\
Semantic object & Persistent metastable region \\
Computation     & Guided entropy navigation \\
\bottomrule
\end{tabular}
\end{center}
\medskip

The resulting interpretation reframes \ambibf{} from an esolang curiosity into
a minimal thermodynamic semantics for mortal computation. Ordinary languages
encourage exact addressing, exact arithmetic, and exact control flow.
\ambibf{} destroys all three, so higher-order structure must emerge
relationally --- making it oddly compatible with RSVP-style field thinking,
semantic manifolds, topological cognition, and distributed attractor systems,
and making the Spherepop layer above it not an awkward workaround but a natural
consequence of the substrate's geometry.

%% ══════════════════════════════════════════════════════════════════════════════
\begin{thebibliography}{99}

\bibitem{barabasi2002}
Barab\'asi, Albert-L\'aszl\'o.
\textit{Linked: The New Science of Networks}.
Cambridge, MA\@: Perseus Publishing, 2002.

\bibitem{barbour2020}
Barbour, Julian, Tim Koslowski, and Flavio Mercati.
\textit{The Janus Point: A New Theory of Time}.
New York: Basic Books, 2020.

\bibitem{benkler2006}
Benkler, Yochai.
\textit{The Wealth of Networks: How Social Production Transforms Markets and Freedom}.
New Haven: Yale University Press, 2006.

\bibitem{brooks1975}
Brooks, Frederick P.
\textit{The Mythical Man-Month: Essays on Software Engineering}.
Reading, MA\@: Addison-Wesley, 1975.

\bibitem{cover2006}
Cover, Thomas M., and Joy A. Thomas.
\textit{Elements of Information Theory}.
2nd ed.
Hoboken, NJ\@: Wiley-Interscience, 2006.

\bibitem{fogel1995}
Fogel, David B.
\textit{Evolutionary Computation: Toward a New Philosophy of Machine Intelligence}.
Piscataway, NJ\@: IEEE Press, 1995.

\bibitem{git2005}
Hamano, Junio C., and Linus Torvalds.
\textit{Git Source Code Management System}.
2005.
\url{https://git-scm.com}

\bibitem{goodfellow2016}
Goodfellow, Ian, Yoshua Bengio, and Aaron Courville.
\textit{Deep Learning}.
Cambridge, MA\@: MIT Press, 2016.

\bibitem{holland1992}
Holland, John H.
\textit{Adaptation in Natural and Artificial Systems}.
Cambridge, MA\@: MIT Press, 1992.

\bibitem{lamport1978}
Lamport, Leslie.
``Time, Clocks, and the Ordering of Events in a Distributed System.''
\textit{Communications of the ACM} 21, no.~7 (1978): 558--565.

\bibitem{kuramoto1984}
Kuramoto, Yoshiki.
\textit{Chemical Oscillations, Waves, and Turbulence}.
Berlin: Springer-Verlag, 1984.

\bibitem{lakatos1976}
Lakatos, Imre.
\textit{Proofs and Refutations}.
Cambridge: Cambridge University Press, 1976.

\bibitem{lessig2001}
Lessig, Lawrence.
\textit{The Future of Ideas}.
New York: Random House, 2001.

\bibitem{neumann1956}
von Neumann, John.
``Probabilistic Logics and the Synthesis of Reliable Organisms from Unreliable
Components.''
In \textit{Automata Studies}, edited by Claude Shannon and John McCarthy,
43--98.
Princeton: Princeton University Press, 1956.

\bibitem{ostrom1990}
Ostrom, Elinor.
\textit{Governing the Commons}.
Cambridge: Cambridge University Press, 1990.

\bibitem{raymond1999}
Raymond, Eric S.
\textit{The Cathedral and the Bazaar}.
Sebastopol, CA\@: O'Reilly Media, 1999.

\bibitem{russell2010}
Russell, Stuart, and Peter Norvig.
\textit{Artificial Intelligence: A Modern Approach}.
3rd ed.
Upper Saddle River, NJ\@: Prentice Hall, 2010.

\bibitem{shannon1948}
Shannon, Claude E.
``A Mathematical Theory of Communication.''
\textit{Bell System Technical Journal} 27, no.~3 (1948): 379--423.

\bibitem{simon1962}
Simon, Herbert A.
``The Architecture of Complexity.''
\textit{Proceedings of the American Philosophical Society} 106, no.~6
(1962): 467--482.

\bibitem{stallman2002}
Stallman, Richard M.
\textit{Free Software, Free Society}.
Boston: GNU Press, 2002.

\bibitem{torvalds2001}
Torvalds, Linus, and David Diamond.
\textit{Just for Fun: The Story of an Accidental Revolutionary}.
New York: HarperBusiness, 2001.

\bibitem{turing1936}
Turing, Alan M.
``On Computable Numbers, with an Application to the Entscheidungsproblem.''
\textit{Proceedings of the London Mathematical Society} 42, no.~2
(1936): 230--265.

\bibitem{vaswani2017}
Vaswani, Ashish, Noam Shazeer, Niki Parmar, Jakob Uszkoreit,
Llion Jones, Aidan N. Gomez, {\L}ukasz Kaiser, and Illia Polosukhin.
``Attention Is All You Need.''
In \textit{Advances in Neural Information Processing Systems~30},
5998--6008. 2017.

\bibitem{wiener1948}
Wiener, Norbert.
\textit{Cybernetics: Or Control and Communication in the Animal and the Machine}.
Cambridge, MA\@: MIT Press, 1948.

\bibitem{wolfram2002}
Wolfram, Stephen.
\textit{A New Kind of Science}.
Champaign, IL\@: Wolfram Media, 2002.

\end{thebibliography}

\end{document}
